%
% File: teach/elmagii/lect-00/lecture-00.tex [plain TeX code]
% Github: https://github.com/elmagii/lect-00/
% Last change: November 4, 2026
%
% Lecture No 0 in the course ``Elektromagnetism II, 1TE626 (2022-2026)'',
% held November 3, 2026, at Uppsala University, Sweden.
%
% Copyright (C) 2022-2026, Fredrik Jonsson, under Gnu General Public
% License (GPL) v3. See the enclosed LICENSE for details.
%
% This program is free software: you can redistribute it and/or modify
% it under the terms of the GNU General Public License as published by
% the Free Software Foundation, either version 3 of the License, or
% (at your option) any later version.
%
% This program is distributed in the hope that it will be useful,
% but WITHOUT ANY WARRANTY; without even the implied warranty of
% MERCHANTABILITY or FITNESS FOR A PARTICULAR PURPOSE.  See the
% GNU General Public License for more details.
%
% You should have received a copy of the GNU General Public License
% along with this program.  If not, see <https://www.gnu.org/licenses/>.
%
\def\coursename{Elektromagnetism II}
\def\coursecode{1TE626}
\def\courseyear{2026}
\def\courserepo{https://github.com/hp35/elmagii/}
\def\lecturenumber{0}
\def\lecturetitle{Grundl{\"a}ggande vektoranalys inom f{\"a}ltteori}
\def\lecturetitlesub{}
\def\lectureauthor{Fredrik Jonsson}
\def\lectureplace{Uppsala Universitet}
\def\lecturedate{3 november 2026}
%-------------------- BEGIN OF LOCAL MACROS --------------------
\input macros/epsf.tex
\input macros/eplain.tex
\input amssym % to get the {\Bbb E} font (strikethrough E)
\font\ninerm=cmr9
\font\tenssbx=cmssbx10
\font\twelvesc=cmcsc10 at 12 truept

%
% Define a macro for the highlighting of text in color (red). In this macro,
% the \newif\ifcolors call creates \ifcolors, \colorstrue, \colorsfalse,
% where \colorstrue turns colors on, \colorsfalse turns them off.
%
\newif\ifcolors
\ifx\colorsenabled\undefined \colorsfalse \else \colorstrue \fi
\ifcolors\input color\fi
\def\hl#1{\ifcolors{\color{red}#1}\else#1\fi\ignorespaces}

\def\ifempty#1{\ifx\relax#1\relax}
\def\initlecture{
  \hsize=150mm\hoffset=4.6mm\vsize=230mm\voffset=7mm
  \topskip=0pt\baselineskip=12pt\parskip=0pt\leftskip=0pt\parindent=15pt
  \ifcolors
    \voffset=-10.2mm\topskip=0pt
  \fi
  \headline={\ifnum\pageno>1\ifodd\pageno\rightheadline\else\leftheadline\fi
    \else\hfill\fi}
  \def\rightheadline{\tenrm{\it F\"orel\"asning \lecturenumber}
    \hfil{\it \coursename, \coursecode\ (\courseyear)}}
  \def\leftheadline{\tenrm{\it \coursename, \coursecode\ (\courseyear)}
    \hfil{\it F\"orel\"asning \lecturenumber}}
  \noindent~\vskip-60pt\hskip-40pt{\epsfbox{macros/UU_logo_color.eps}}
  \vskip-42pt\hfill\vbox{
      \hbox{{\it \coursename, \coursecode\ (\courseyear)}}
      \hbox{{\it Lecture Notes, \lectureauthor}}
      \hbox{{\it Document Revision \today}}
      \hbox{{\it \courserepo}}}\vskip 36pt
  \centerline{\twelvesc F\"orel\"asning \lecturenumber}
  \vskip 24pt\noindent
  \centerline{\twelvesc\lecturetitle}
  \expandafter\ifempty\expandafter{\lecturetitlesub}%
    \else\centerline{\twelvesc\lecturetitlesub}\fi
  \bigskip
  \centerline{\lectureauthor, \lectureplace, \lecturedate}
  \vskip24pt}
\def\section #1 {\medskip\goodbreak\noindent{\tenssbx #1}
  \par\nobreak\smallskip\noindent}
\def\subsection #1 {\medskip\goodbreak\noindent{\it #1}
  \par\nobreak\smallskip\noindent}
\def\iint{\mathop{\int\kern-8pt\int}}
\def\iiint{\mathop{\int\kern-8pt\int\kern-8pt\int}}
\def\oiint{\mathop{\int\kern-8pt\int\kern-13.2pt{\bigcirc}}}
\def\sgn{\mathop{\rm sgn}\nolimits} % sign
\def\Re{\mathop{\rm Re}\nolimits}   % real part
\def\Im{\mathop{\rm Im}\nolimits}   % imaginary part
\def\Tr{\mathop{\rm Tr}\nolimits}   % quantum mechanical trace
\def\eqq{\mathop{\vbox{\hbox{\hskip2pt?}\vskip-6pt\hbox{=}}}}
\def\quote#1{\par\leftskip=36pt\rightskip=36pt\smallskip\noindent#1\par
  \leftskip=0pt\rightskip=0pt\smallskip}
\long\def\plan#1{\par\leftskip=36pt\rightskip=36pt\bigskip
  \noindent{\it Sammanfattning}\smallskip
  \noindent{\it #1}\par\leftskip=0pt\rightskip=0pt}
\def\threepointsummary#1#2#3{\par\leftskip=36pt\rightskip=36pt\bigskip
  \noindent{\it Tre h{\aa}llpunkter i f{\"o}rel{\"a}sningen}\smallskip
  \leftskip=48pt\rightskip=36pt\hangindent=20pt
  \noindent{\it\hbox to 20pt{1. }#1}\smallskip
  \leftskip=48pt\rightskip=36pt\hangindent=20pt
  \noindent{\it\hbox to 20pt{2. }#2}\smallskip
  \leftskip=48pt\rightskip=36pt\hangindent=20pt
  \noindent{\it\hbox to 20pt{3. }#3}\par
  \leftskip=0pt\rightskip=0pt\vfill\eject}
\def\epsfig#1{\bigskip\centerline{\epsfbox{#1}}\medskip}
\def\captionwide{\advance\leftskip by 60pt
  \advance\rightskip by 60pt}
\newif\ifshowindex
\showindextrue  % Use \showindextrue and \showindexfalse to enable/disable index
\def\index{\ifshowindex\vfill\eject\section{Index} \readindexfile{i}\fi}
%--------------------- END OF LOCAL MACROS ---------------------

\initlecture

\plan{Denna f{\"o}rel{\"a}sning {\"a}r bist{\aa}r med ett extramaterial som
avhandlar de grundl{\"a}ggande begrepp inom vektoranalys som vi frekvent kommer
att anv{\"a}nda oss av i kommande f{\"o}rel{\"a}sningar.
Denna extraf{\"o}rel{\"a}sning h{\aa}lls i n{\"a}ra anslutning till den
f{\"o}rsta f{\"o}rel{\"a}sningen i kursen, och avser att fungera som
rekapitulation f{\"o}r den teknolog som k{\"a}nner sig lite ringrostig
i sin vektoranalys.

Begreppen f{\"a}lt, gradient, divergens och rotation g{\aa}s igenom med
praktiska exempel h{\"a}mtade fr{\aa}n elektrostatik och magnetostatik.
Vi inleder med en rekapitulation kring begreppet delta-distribution i en
och tre dimensioner, f{\"o}ljt av en genomg{\aa}ng av de grundl{\"a}ggande
definitionerna av gradient $\nabla\phi$, divergens $\nabla\cdot{\bf a}$ och
rotation $\nabla\times{\bf b}$.
De grundl{\"a}ggande identiteterna f{\"o}r divergensen av en rotation och
rotationen av en gradient g{\aa}s igenom, d{\aa} dessa {\"a}r v{\"a}sentliga
f{\"o}r f{\"o}rst{\aa}elsen av hur vi rent matematiskt kan motivera existensen
av vektorpotentialen och den skal{\"a}ra potentialen.
F{\"o}rel{\"a}sningen avslutas med att illustrera vad rotationen av en rotation
($\nabla\times\nabla\times$) inneb{\"a}r, d{\aa} detta {\"a}r centralt f{\"o}r
formuleringen av den elektromagnetiska v{\aa}gekvationen.}

\threepointsummary{%
  Diracs delta-distribution $\delta({\bf x})$, som kan anv{\"a}ndas som en
  s{\"a}tt att f{\aa} med punkt, linje- och yt-laddningar i en generell
  beskrivning i laddningst{\"a}theten $\rho({\bf x})$ eller
  str{\"o}mt{\"a}theten ${\bf J}({\bf x})$.
  Exempel:
    (A) {\it Punktladdning} $q$ i k{\"a}llpunkten ${\bf x}'$, be\-skriven av
        volymladdningst{\"a}theten $\rho({\bf x})=q\delta({\bf x}-{\bf x}')$,
    (B) {\it Linjeladdning} med linjeladdningst{\"a}thet $\lambda$ l{\"a}ngs
        $z$-axeln och genom $(x',y')$, be\-skriven av $\rho({\bf x})=\lambda
        \delta(x-x')\delta(y-y')$, samt
    (C) {\it Ytladdning} med ytladdningst{\"a}theten $\sigma$ i planet $z=0$,
    beskriven av $\rho({\bf x})=\sigma\delta(z-z')$.
}{%
  Tolkningen av gradienten $\nabla\phi$, divergensen $\nabla\cdot{\bf a}$
  och rotationen $\nabla\times{\bf b}$, samt vad begreppet konservativt
  f{\"a}lt inneb{\"a}r.
}{%
  Tolkningen av vektoridentiteterna
  $$
    \eqalign{
      &\nabla\times(\nabla\phi)=0,\cr
      &\nabla\cdot(\nabla\times{\bf a})={\bf 0},\cr
      &\nabla\times(\nabla\times{\bf a})
         =\nabla(\nabla\cdot{\bf a})-\nabla^2{\bf a},\cr
    }
  $$
  samt hur dessa direkt ger oss ett antal fundamentala relationer inom
  klassisk elektromagnetism.
}
%----------------------- END OF PREAMBLE -----------------------

\section{Introduktion}
Inom elektromagnetism st{\"o}ter vi p{\aa} n{\aa}gra centrala storheter, som
elektrisk laddning, elektriskt f{\"a}lt, magnetiskt f{\"a}lt och elektrisk
potential. Dessa storheter beskriver olika aspekter av hur elektriska och
magnetiska fenomen fungerar, och formulerar de variabler vi har att arbeta
med.\sidx{Elektromagnetism}

Vi kan dock fr{\aa}ga oss hur vi p{\aa} b{\"a}sta s{\"a}tt egentligen beskriver
ett f{\"a}lt som varierar fr{\aa}n punkt till punkt? Hur kan vi till exempel
matematiskt skilja mellan ett f{\"a}lt som pekar ut fr{\aa}n en punkt och ett
som cirkulerar runt den? Framf{\"o}r allt, vilka grundl{\"a}ggande matematiska
operationer kan vi utnyttja som v{\aa}rt ``spr{\aa}k'' f{\"o}r att formulera
teorin f{\"o}r de fysikaliska lagar som g{\"a}ller f{\"o}r elektromagnetiska
f{\"a}lt och deras utbredning i rum och tid? Det {\"a}r h{\"a}r vektoranalysen
kommer till v{\aa}r assistans.\sidx{Vektoranalys}[Introduktion till]

I denna extra f{\"o}rel{\"a}sning kommer vi att i huvudsak introducera fem
grundl{\"a}ggande begrepp, som successivt genom kursen kommer att s{\"a}ttas
samman och kombineras:
\quote{%
  \item{1.}{{\hl{{\it Delta-funktionen} $\delta({\bf x})$}}, {\"a}ven mer
            formellt kallad {\hl{\it Diracs delta-distribution}}, {\"a}r en
            distribution\numberedfootnote{I denna kurs kan vi lite l{\"o}st
              s{\"a}ga att en {\it distribution} {\"a}r en generaliserad
              funktion som {\"a}r definierad av hur den beter sig n{\"a}r
              den {\"a}r en del av en integrand. Vi kommer strax att se
              exempel p{\aa} till{\"a}mpningen av denna definition.}
            som vi kan anv{\"a}nda f{\"o}r ett diskret urval.
            \sidx{Diracs delta-distribution $\delta({\bf x})$}}
  \item{1.}{{\hl{{\it Gradienten}
            ${\rm grad}(\phi)\equiv\nabla\phi$}}
            beskriver hur en skal{\"a}r storhet $\phi$ f{\"o}r{\"a}ndras
            {\it spatialt}.\sidx{Gradient}\numberedfootnote{%
            $\hbox{``Spatialt''}=\hbox{i rummet}$;
            $\hbox{``Temporalt''}=\hbox{i tiden}$.}}
  \item{2.}{{\hl{{\it Divergensen}
            ${\rm div}({\bf a})\equiv\nabla\cdot{\bf a}$}}
            av ett vektorf{\"a}lt ${\bf a}$ beskriver ett vektorf{\"a}lts
            lokala nettofl{\"o}de ut fr{\aa}n en punkt.
            \sidx{Divergens}[Definition]}
  \item{3.}{{\hl{{\it Rotationen}
            ${\rm rot}({\bf a})\equiv\nabla\times{\bf a}$}}
            beskriver ett vektorf{\"a}lts {\it lokala cirkulation}.
            \sidx{Rotation}[Definition]}
  \item{4.}{{\hl{Ett f{\"a}lt ${\bf F}$ {\"a}r {\it konservativt} om dess
            linjeintegral $\int^B_A{\bf F}\cdot d{\bf l}$ mellan tv{\aa}
            punkter inte beror p{\aa} vilken v{\"a}g vi v{\"a}ljer f{\"o}r
            integralen.}}
            \sidx{Konservativt f{\"a}lt}[Definition]
            \sidx{F{\"a}lt}[Konservativt]}
}
\noindent
I denna extra f{\"o}rel{\"a}sning kommer vi att successivt f{\"o}rs{\"o}ka
knyta de introducerade begreppen inom vektoranalys till v{\aa}rt specifika
{\"a}mne, elektromagnetism, som praktiskt taget {\"a}r det f{\"a}lt inom
fysiken som vektoranalysen utvecklades till att ge ett matematiskt spr{\aa}k
f{\"o}r. M{\aa}let {\"a}r inte bara att l{\"a}ra sig r{\"a}kna med operatorn
$\nabla$, utan att f{\"o}rst{\aa} vad ber{\"a}kningarna betyder fysikaliskt.

De exempelfall som vi h{\"a}r g{\aa}r igenom {\"a}r sj{\"a}lvfallet mycket
f{\"o}renklade, f{\"o}r att illustrera de koncept vi rekapitulerar som
delta-funktioner och differentialoperatorerna grad, div och rot.
Det {\"a}r dock mycket enkelt att generalisera dessa till godtyckliga
distributioner {\"o}ver kurvor, ytor och volymer. Med andra ord skall
materialet i f{\"o}rel{\"a}sningen tas som en bas f{\"o}r vidare utveckling
av analysen och f{\"a}rdigheter inom probleml{\"o}sande inom elektromagnetism.

\section{\hl{Vad {\"a}r ett f{\"a}lt?}}
\sidx{F{\"a}lt}[Definition]
Vi b{\"o}rjar med en enkel fr{\aa}ga: {\it Vad menar vi egentligen n{\"a}r vi
s{\"a}ger att en storhet {\"a}r ett f{\"a}lt?}
\quote{%
  {\hl{{\bf Definition.} Ett f{\"a}lt {\"a}r en storhet som tilldelar ett
  v{\"a}rde till varje punkt i ett omr{\aa}de av rummet, och eventuellt
  ocks{\aa} till varje tidpunkt.}}
}
\noindent
Rent matematiskt beskriver vi ett f{\"a}lt som en funktion av
positionen\numberedfootnote{Genomg{\aa}ende i denna kurs kommer vi att
  anv{\"a}nda ${\bf x}$ f{\"o}r {\it ortsvektorn} till en punkt i rummet,
  och i de fall vi har k{\"a}llpunkter, exempelvis f{\"o}r laddningar
  eller str{\"o}mmar, s{\aa} kommer vi f{\"o}r dessa att anv{\"a}nda
  primmade koordinater ${\bf x}'$, s{\aa} att vektorn fr{\aa}n k{\"a}lla
  till observationspunkt {\"a}r ${\bf x}-{\bf x}'$}
${\bf x}$ och tiden $t$.

Det finns tv{\aa} grundl{\"a}ggande typer av f{\"a}lt som vi kommer att
st{\"o}ta p{\aa} i denna kurs:
\quote{%
  \item{{\hl{1.}}}{{\hl{Skal{\"a}rf{\"a}lt}} tilldelar varje punkt ett enda
    tal, med $\phi=\phi({\bf x},t)$. Typiska exempel {\"a}r temperaturen i
    ett rum, lufttrycket i atmosf{\"a}ren och den elektriska skal{\"a}ra
    potentialen.}
  \item{{\hl{2.}}}{{\hl{Vektorf{\"a}lt}} tilldelar ist{\"a}llet varje punkt
    en vektor, ${\bf a}={\bf a}({\bf x},t)$. Typiska exempel {\"a}r
    vindhastigheten i atmosf{\"a}ren, hastighetsf{\"a}ltet i en
    v{\"a}tska, det elektriska f{\"a}ltet ${\bf E}$ och det magnetiska
    f{\"a}ltet ${\bf B}$.}
}
\noindent
Ett vektorf{\"a}lt inneh{\aa}ller allts{\aa} information om b{\aa}de
{\it riktning och storlek} i varje punkt ${\bf x}$.

Vi kan t{\"a}nka oss att vi placerar en liten pil i varje punkt i rummet.
Pilens riktning och l{\"a}ngd beskriver vektorf{\"a}ltet d{\"a}r.
Om vi i st{\"a}llet anger en temperatur med ett tal i varje punkt f{\aa}r
vi ett skal{\"a}rf{\"a}lt.

\subsection{Varifr{\aa}n kommer ordet ``f{\"a}lt''?}
Ordet ``f{\"a}lt'' {\"a}r kort och gott bara en rumslig metafor. I v{\aa}rt
vardagsspr{\aa}k {\"a}r ett f{\"a}lt ett omr{\aa}de som str{\"a}cker sig
{\"o}ver en viss yta. Inom fysiken har begreppet f{\aa}tt en mer precis
betydelse: n{\aa}got som {\it har ett v{\"a}rde {\"o}verallt i ett omr{\aa}de},
och inte bara p{\aa} en enskild plats.

Historiskt blev f{\"a}ltbegreppet s{\"a}rskilt viktigt under 1800-talet, bland
annat genom Michael Faradays\numberedfootnote{Vi kommer frekvent att
  {\aa}terkomma till Michael Faraday under kommand f{\"o}rel{\"a}sningar i
  denna kurs.} \sidx{Faraday, Michael (1791--1867)}
arbete med elektriska och magnetiska fenomen. Faraday anv{\"a}nde
f{\"o}rest{\"a}llningen om kraftlinjer f{\"o}r att beskriva hur elektriska
och magnetiska krafter f{\"o}rdelar sig i rummet.
James Clerk Maxwell\sidx{Maxwell, James Clerk (1831--1879)}, som vi kommer att
{\aa}terkomma till i F{\"o}rel{\"a}sning~9 d{\aa} vi introducerar den s{\aa}
kallade {\it f{\"o}rskjutningsstr{\"o}mmen}, utvecklade senare en matematisk
teori d{\"a}r f{\"a}lten blev centrala fysikaliska storheter.

En viktig po{\"a}ng {\"a}r att ett elektromagnetiskt f{\"a}lt inte bara {\"a}r
ett s{\"a}tt att rita kraftlinjer. F{\"a}ltet {\"a}r en fysikalisk storhet som
kan b{\"a}ra energi och r{\"o}relsem{\"a}ngd och som kan fortplanta sig genom
rummet i form av elektromagnetiska v{\aa}gor.

\subsection{Elektriskt och magnetiskt f{\"a}lt}
I denna kurs kommer vi prim{\"a}rt att arbeta med tre vektorf{\"a}lt:
\quote{%
  \item{1.}{{\it Det elektriska f{\"a}ltet} ${\bf E}$, som kan definieras genom
    den elektriska kraft ${\bf F}=q{\bf E}$ som genom Coulomb-v{\"a}xelverkan
    verkar p{\aa} en testladdning $q$.
    I denna kurs betraktar vi det elektriska f{\"a}ltet som definierat genom
    ${\bf E}={\bf F}/q$, d{\"a}r vi t{\"a}nker oss att testladdningen $q$
    {\"a}r tillr{\"a}ckligt liten f{\"o}r att inte p{\aa}verka det ursprungliga
    f{\"a}ltet p{\aa} n{\aa}got v{\"a}sentligt s{\"a}tt.
    Det elektriska f{\"a}ltet beskriver kraften per laddningsenhet.
    Enheten {\"a}r newton per coulomb, vilket {\"a}r ekvivalent med volt
    per meter.\sidx{Elektrisk f{\"a}ltstyrka ${\bf E}$}[Dynamisk, tidsberoende]}
  \item{2.}{{\it Det magnetiska f{\"a}ltet} ${\bf B}$ kan p{\aa} motsvarande
    s{\"a}tt kopplas till den magnetiska delen av Lorentz-kraften p{\aa}
    en laddning i r{\"o}relse med hastigneten ${\bf v}$ genom
    ${\bf F}_{\rm m}=q\,{\bf v}\times{\bf B}$, vilken {\"a}r vinkelr{\"a}t
    mot b{\aa}de hastigheten och magnetf{\"a}ltet.\numberedfootnote{Om
      vi skall vara riktigt petiga, s{\aa} skall ${\bf B}$ betecknas
      {\it magnetisk fl{\"o}dest{\"a}thet}.}
    {\it Till skillnad fr{\aa}n det elektriska f{\"a}ltet kan det magnetiska
    f{\"a}ltet inte generellt best{\"a}mmas genom att dividera en kraft med
    laddningen utan att ocks{\aa} ta h{\"a}nsyn till laddningens hastighet
    och r{\"o}relseriktning.}
    \sidx{Lorentz-kraften ${\bf F}_{\rm m}=q\,{\bf v}\times{\bf B}$}}
}

\subsection{Skal{\"a}r potential och vektorpotential}
Inom elektromagnetism anv{\"a}nder vi {\"a}ven tv{\aa} andra f{\"a}lt, kallade
den {\it skal{\"a}ra potentialen}\numberedfootnote{Ofta lite slappt kallad
  ``den elektriska skal{\"a}ra potentialen''} $\phi({\bf x},t)$ och
{\it vektorpotentialen}\numberedfootnote{{\AA}terigen, ofta lite slappt kallad
  ``den magnetiska vektorpotentialen''}
${\bf A}({\bf x},t)$.

I {\it elektrostatiska} problem kan vi det uttrycka elektriska f{\"a}ltet som
$$
  {\bf E}=-\nabla\phi,
$$
d{\"a}r den {\it skal{\"a}ra potentialen} $\phi$ {\"a}r ett {\it skal{\"a}rt
f{\"a}lt} d{\"a}r varje punkt har ett skal{\"a}rt v{\"a}rde f{\"o}r potentialen,
m{\"a}tt i enheten volt. Gradienten $\nabla\phi$ beskriver hur potentialen
f{\"o}r{\"a}ndras i rummet, och minustecknet inneb{\"a}r att det
{\it elektriska f{\"a}ltet pekar i riktningen f{\"o}r den snabbaste minskning
av potentialen}.

F{\"o}r magnetf{\"a}ltet, eller mer korrekt den {\it magnetiska
fl{\"o}dest{\"a}theten}, har vi motsvarande relation som
$$
  {\bf B}=\nabla\times{\bf A},
$$
d{\"a}r den underliggande {\it vektorpotentialen} ${\bf A}$ {\"a}r ett
{\it vektorf{\"a}lt}, och d{\"a}r rotationen av detta vektorf{\"a}lt ger
det magnetiska f{\"a}ltet ${\bf B}$.
Vektorpotentialen ${\bf A}$ ligger som grund f{\"o}r beskrivningen ``under
huven'' av statiska magnetiska f{\"a}lt\numberedfootnote{Som vi kommer in
  p{\aa} i detalj i F{\"o}rel{\"a}sning~5 och~7.}~${\bf B}$ och f{\"o}r
tidsveroende elektriska f{\"a}lt\numberedfootnote{Som vi kommer in p{\aa}
  i detalj i F{\"o}rel{\"a}sning~10 och~11.}~${\bf E}$.
Vektorpotentialen {\"a}r dock av den naturen att vi inte kan m{\"a}ta den med
ett m{\"a}tinstrument, utan att vi enbart kan m{\"a}ta dess {\it observabler}
i form av~${\bf E}$ och~${\bf B}$. Trots detta har dock vektorpotentialen en
fysikalisk enhet, som volt-sekunder per meter (${\rm V}\cdot{\rm s}/{\rm m}$).
\sidx{Vektorpotential ${\bf A}$}
\sidx{Elektrisk f{\"a}ltstyrka ${\bf E}$}[Dynamisk, tidsberoende]
\sidx{Magnetisk fl{\"o}dest{\"a}thet ${\bf B}$}[Statisk]

Vi kommer i kommande f{\"o}rel{\"a}sningar att frekvent n{\"a}mna ursprunget
till hur det kommer sig att vektorpotentialen existerar som det f{\"a}lt som
ligger bakom magnetiska f{\"a}lt, men l{\aa}t oss trots detta redan nu
klarg{\"o}ra hur det kommer sig att vektorpotentialen och det magnetiska
f{\"a}ltet {\"a}r kopplade p{\aa} detta s{\"a}tt.

Maxwells ekvationer ger oss bland annat Gauss lag f{\"o}r magnetiska
fl{\"o}dest{\"a}theten som $\nabla\cdot{\bf B}=0$, och i ett ``vanligt
omr{\aa}de'' av rummet utan singulariteter\numberedfootnote{Med
  ``singulariteter'' menar vi h{\"a}r skarpa v{\"a}ggar, diskreta
  {\"o}verg{\aa}ngar mellan media med olika magnetiska permeabiliteter
  i form av randv{\"a}rden, och s{\aa} vidare.}
inneb{\"a}r detta att magnetf{\"a}ltet {\"a}r divergensfritt. Under l{\"a}mpliga
matematiska f{\"o}ruts{\"a}ttningar kan ett s{\aa}dant f{\"a}lt representeras
som rotationen av en vektorpotential.\numberedfootnote{Notera att de
  elektromagnetiska potentialerna $\phi$ och ${\bf A}$ inte {\"a}r
  unika. Vi kan exempelvis {\"a}ndra vektorpotentialen enligt
  ${\bf A}'={\bf A}+\nabla\chi$ f{\"o}r en godtyckligt vald funktion
  $\chi({\bf x})$ utan att {\"a}ndra magnetf{\"a}ltet ${\bf B}$, eftersom
  rotationen av en gradient {\"a}r noll. Detta kallas {\it gaugefrihet}
  under {\it gaugetransformation}, och {\"a}r n{\aa}got som vi {\"o}versiktligt
  kommer att g{\aa} in p{\aa} i F{\"o}rel{\"a}sning~11.}
I den h{\"a}r f{\"o}rel{\"a}sningen kommer vi att se att gradient, divergens
och rotation {\"a}r de tre centrala verktygen som g{\"o}r dessa samband
begripliga.

\section{Diracs delta-distribution och punktformiga k{\"a}llor}
Innan vi b{\"o}rjar beskriva elektriska och magnetiska f{\"a}lt med gradient,
divergens och rotation, som alla {\"a}r operatorer som rent principiellt
arbetar p{\aa} ``sl{\"a}ta'' funktioner och f{\"a}lt, s{\aa} beh{\"o}ver vi
ett matematiskt verktyg f{\"o}r att beskriva n{\aa}got som {\"a}r mycket
koncentrerat till en viss punkt i rummet. Dessa ``koncentrat'' kan exempelvis
vara en diskret punktladdning,\numberedfootnote{I denna
  f{\"o}rel{\"a}sningsserie betecknar vi ofta punktladdningen
  som en idealiserad ``punktk{\"a}lla'', Vi {\aa}terkommer strax
  till varf{\"o}r en punktladdning kan ses som en k{\"a}lla inom
  elektromagnetiska problem.}
ett system av diskreta punktladdningar, eller en rent av en impuls av ett
f{\"a}lt som verkar under ett mycket kort tidsintervall. Detta matematiska
verktyg {\"a}r den s{\aa} kallade {\it Diracs delta-distribution}, ofta kort och
gott kallad ``delta-puls''.\numberedfootnote{Vi kommer i denna kurs att
  f{\"o}rs{\"o}ka h{\aa}lla oss till att genomg{\aa}ende s{\"a}ga
  {\it delta-distribution} och inte ``delta-funktion''; det {\"a}r
  dock h{\"o}gst sannolikt att f{\"o}rel{\"a}saren vid ett flertal
  tillf{\"a}llen oavsiktligt kommer att s{\"a}ga just ``delta-funktion'',
  och det {\"a}r d{\aa} helt okej att korrigera under p{\aa}g{\aa}ende
  f{\"o}rel{\"a}sning!}\sidx{Diracs delta-distribution $\delta({\bf x})$}

Denna distribution {\"a}r s{\"a}rskilt viktig inom elektromagnetismen, eftersom
vi ofta vill beskriva laddningar och str{\"o}mmar som {\"a}r koncentrerade till
mycket sm{\aa} omr{\aa}den, men samtidigt vill anv{\"a}nda de generella
differentialekvationer som g{\"a}ller i hela rummet och som alltsom oftast
{\"a}r definierade via integraler.
Genom att anv{\"a}nda Diracs delta-distribution ges vi m{\"o}jlighet att
inkludera just diskreta laddningar definierade i punkter, linjer eller ytor,
trots att differentialoperatorerna opererar p{\aa} kontinuerliga
distributioner, som vi nu skall se.
\vfill\eject

\subsection{Delta-distributionen i en dimension -- En mycket smal puls}
Vi f{\"o}rest{\"a}ller oss en puls l{\"a}ngs $x$-axeln, exempelvis en
rektangul{\"a}r puls som {\"a}r noll {\"o}verallt utom i ett litet intervall
kring $x=0$.
Om vi g{\"o}r pulsen allt smalare, men samtidigt allt h{\"o}gre, s{\aa} att
arean under kurvan f{\"o}rblir lika stor, och l{\aa}ter pulsebredden vara
$\Delta x=\epsilon$ och h{\"o}jden $1/\epsilon$, s{\aa} blir naturligtvis
pulsens area trivialt
$$
  \hbox{area}=\epsilon\cdot{{1}\over{\epsilon}}=1.
$$
D{\aa} vi l{\aa}ter pulsbredden $\epsilon$ g{\aa} mot noll och h{\"o}jden
$1/\epsilon$ d{\"a}rmed mot o{\"a}ndligheten, s{\aa} f{\aa}r vi en idealiserad
gr{\"a}ns d{\"a}r pulsen {\it praktiskt taget endast existerar i en punkt}, och
det {\"a}r denna {\it distribution} som vi betecknar som {\it delta-funktionen}
$\delta(x)$ i en dimension.
\medskip
\epsfig{figs/dirac-pulse.1}\noindent
\medskip
\noindent
Vi kan sammanfatta delta-distributionens egenskaper med att den {\"a}r noll
{\"o}verallt f{\"o}rutom i origo,
$$
  \delta(x)=0,\quad\hbox{f{\"o}r $x\ne0$},
$$
samt att varje integral som t{\"a}cker in origo kommer att ge ett enhetsbidrag,
d{\aa} pulsen {\"a}r normaliserad som
$$
  \int^{\infty}_{-\infty}\delta(x)\,dx=1.
$$
F{\"o}r att sammanfatta delta-distributionen s{\aa} {\"a}r den allts{\aa}
koncentrerad till $x=0$, samt att den {\"a}r noll {\"o}verallt utom i en
enda punkt.\sidx{Diracs delta-distribution $\delta({\bf x})$}[En dimension]

H{\"a}r m{\aa}ste vi dock g{\"o}ra en viktig matematisk reservation. Diracs
delta-distribution {\"a}r definitivt inte en ``vanlig funktion'' i samma mening
som exempelvis $\sin(x)$ eller $\exp(-x^2)$. Den {\"a}r som namnet s{\"a}ger en
{\it distribution}, definierad genom hur den verkar n{\"a}r den anv{\"a}nds
i integranden i en integral.
F{\"o}rest{\"a}llningen om en o{\"a}ndligt h{\"o}g och o{\"a}ndligt smal puls
{\"a}r h{\"a}r en anv{\"a}ndbar bild, men inte en fullst{\"a}ndig matematisk
definition, eftersom {\"a}ven en ``kantig'' distribution st{\"a}ller till det
i vissa aspekter av differentierbarhet.
I denna f{\"o}rel{\"a}sningsserie beh{\"o}ver vi dock inte bry oss om dessa
aspekter, d{\aa} vi enbart kommer att anv{\"a}nda oss av delta-distributionen
just som ing{\aa}ende som integrand.
\vfill\eject

\subsection{Den viktigaste egenskapen hos delta-distributionen {\"a}r
  urvalsegenskapen}
Den verkliga styrkan hos Diracs delta-funktion ligger i att den kan selektera
(``plocka ut'') v{\"a}rdet p{\aa} en annan funktion i en best{\"a}md punkt,
s{\"a}g vid $x=a$ i ett endimensionellt fall. Om $f(x)$ {\"a}r en
tillr{\"a}ckligt v{\"a}lbeteende funktion s{\aa} g{\"a}ller det att integralen
av funktionen, n{\"a}r den multipliceras med delta-distributionen med en
offset $a$, till att ha argumentet $(x-a)$, blir
$$
  \int^{\infty}_{-\infty}f(x)\delta(x-a)\,dx=f(a),
$$
vilket kan s{\"a}gas just {\it selektera} v{\"a}rdet $f(a)$ via anv{\"a}ndandet
av delta-distributionen. Detta fungerar genom att delta-distributionen, i
gr{\"a}nsen d{\aa} dess bredd $\epsilon\to0$, {\"a}r koncentrerad till just
punkten $x=a$, och integralen bidrar d{\"a}rf{\"o}r endast v{\"a}rdet av $f(x)$
i just punkten $x=a$.
Det {\"a}r denna egenskap, snarare {\"a}n den o{\"a}ndliga h{\"o}jden, som
{\"a}r central f{\"o}r alla till{\"a}mpningar.\numberedfootnote{Tv{\aa}
  exempel p{\aa} till{\"a}mpningen av denna selektion {\"a}r
  $$
    \int^{\infty}_{-\infty}x^2\delta(x-3)\,dx=9,
    \quad\hbox{och}\quad
    \int^{\infty}_{-\infty}\sin(x)\delta(x-\pi/2)\,dx=1.
  $$}
Delta-distributionen beh{\"o}ver allts{\aa} inte sj{\"a}lv vara en fysisk puls.
Den fungerar som ett generellt matematiskt instrument f{\"o}r att representera
en lokaliserad k{\"a}lla.
\sidx{Diracs delta-distribution $\delta({\bf x})$}[Selektionsegenskap]

\subsection{Delta-distributionen med en annan total styrka}
Antag att vi vill beskriva en puls med en total integral $q$, till exempel
m{\aa}ttet p{\aa} en elektrisk punktladdning, i st{\"a}llet f{\"o}r $1$.
I detta fall kan vi sj{\"a}lvfallet ist{\"a}llet anv{\"a}nda
$$
  f(x)=q\delta(x-a)\qquad\Rightarrow\qquad\int^{\infty}_{-\infty}f(x)\,dx=q,
$$
d{\"a}r faktorn $q$ anger pulsens {\it totala styrka}, och d{\"a}r
delta-distributionen i sig anger var den {\"a}r {\it lokaliserad}.
Detta {\"a}r en metod som vi direkt i F{\"o}rel{\"a}sning~1 kommer att
anv{\"a}nda oss av f{\"o}r att beskriva punktladdningar.

\subsection{Fr{\aa}n en dimension till tre dimensioner}
I en dimension lokaliserar, eller snarare {\it koncentrerar}, vi
delta-distributionen till en punkt p{\aa} en {\it linje}. I tre dimensioner,
d{\"a}r vi ist{\"a}llet beh{\"o}ver lokalisera n{\aa}got till en enda punkt i
{\it rummet}, exempelvis en k{\"a}llpunkt ${\bf x}'=(x',y',z')$,
kan vi anv{\"a}nda den tredimensionella delta-distributionen, vilken uttryckt
i kartesiska koordinater lyder\numberedfootnote{Alltsom oftast kommer vi att
  i denna serie av f{\"o}rel{\"a}sningar att skippa ``${}^{(3)}$'' i notationen,
  d{\aa} det utifr{\aa}n argumentet {\"a}nd{\aa} kommer att vara helt klart
  huruvida distributionen har ett endimensionellt (skal{\"a}rt) eller
  tredimensionellt (vektor) argument.}
  \sidx{Diracs delta-distribution $\delta({\bf x})$}[Tre dimensioner]
$$
  \delta^{(3)}({\bf x}-{\bf x}')=\delta(x-x')\delta(y-y')\delta(z-z').
$$
Denna distribution {\"a}r koncentrerad till den tredimensionella punkten
${\bf x}'$, och helt analogt med det endimensionella fallet g{\"a}ller
selektionssegenskapen p{\aa} samma s{\"a}tt, med enda skillnaden att den
tredimensionella delta-funktionen anv{\"a}nds i en {\it volymintegral} enligt
$$
  \iiint_{{\Bbb R}^3}f({\bf x})\delta^{(3)}({\bf x}-{\bf x}')\,dV = f({\bf x}')
$$
d{\"a}r $dV=dx\,dy\,dz$ {\"a}r volymelementet i koordinatsystemet
${\bf x}=(x,y,z)$.\sidx{Volymintegral}[Av Diracs delta-distribution]
Liksom f{\"o}r det skal{\"a}ra, endimensionella fallet {\"a}r
delta-distributionen {\"a}ven i tre dimensioner normaliserad, med
$$
  \int_{{\Bbb R}^3}\delta^{(3)}({\bf x}-{\bf x}')\,dV=1.
$$
En viktig detalj {\"a}r att delta-funktionen i tre dimensioner har den
fysikaliska dimensionen av {\it invers volym} ($1/{\rm m}^3$), om koordinaterna
m{\"a}ts i meter. I en dimension har den p{\aa} motsvarande s{\"a}tt den
fysikaliska dimensionen av {\it invers l{\"a}ngd} ($1/{\rm m}$).
Detta f{\"o}ljer helt naturligt av att integralen av delta-funktionen {\"o}ver
det omr{\aa}de d{\"a}r den {\"a}r koncentrerad skall vara dimensionsl{\"o}s,
och dessutom ut{\"o}ver detta normaliserad- (Volymselementet $dV$ har ju den
fysikaliska dimensionen av just volym (${\rm m}^3$.)

\subsection{En punktladdning}
Antag nu att vi har placerat en elektrisk punktladdning $q$ i k{\"a}llpunkten
${\bf x}'$. Vanligtvis beskriver vi inom ramverket f{\"o}r elektromagnetisk
f{\"a}ltteori en laddningsf{\"o}rdelning i rummet med den kontinuerliga
laddningst{\"a}theten $\rho({\bf x})$, med enheten coulomb per kubikmeter
(${\rm C}/{\rm m}^3$). Laddningen i ett infinitesimalt litet volymelement
$dV'$ vid k{\"a}llpunkten ${\bf x}'$ {\"a}r d{\aa} $dq=\rho({\bf x}')\,dV$.

Om vi ist{\"a}llet f{\"o}r en kontinuerlig laddningst{\"a}thet vill beskriva
en idealiserad punktladdning, som matematiskt verkligen {\"a}r placerad i en
``noll-dimensionell punkt'' ${\bf x}'$, s{\aa} kan vi anv{\"a}nda
delta-distributionen f{\"o}r att beskriva denna punktladdning med hj{\"a}lp
av {\it distributionen av laddningst{\"a}thet i rummet} som
$$
  \rho({\bf x})=q\delta^{(3)}({\bf x}-{\bf x}').
$$
Detta kan f{\"o}rst verka m{\"a}rkligt. Hur kan en laddningst{\"a}thet, via
delta-distributionen $\delta^{(3)}({\bf x}-{\bf x}')$, vara noll n{\"a}stan
{\"o}verallt, men {\"a}nd{\aa} ge en total laddning $q$?

Svaret p{\aa} denna paradox\numberedfootnote{En {\it paradox} {\"a}r ett p{\aa}st{\aa}ende,
  en tanke eller en situation som vid f{\"o}rsta anblicken verkar mots{\"a}gelsefull
  eller orimlig, men som {\"a}nd{\aa} kan rymma en djupare sanning.}
ges av att vi definierat punktladdningen via en distribution, som i sig saknar
egentlig analytisk mening fram till den punkt d{\"a}r den s{\"a}tts i sitt
sammanhang i form av en integral, i detta fall en integral {\"o}ver alla
rumskordinater som ing{\aa}r i dom{\"a}nen f{\"o}r problembeskrivningen.

Vi ber{\"a}knar den totala laddningen $q_{\rm tot}$ genom att integrera
laddningst{\"a}theten $\rho({\bf x})$ {\"o}ver hela rummet i
fr{\aa}ga,\numberedfootnote{Att ``integrera {\"o}ver hela rummet'' {\"a}r
  h{\"a}r som vanligt en omskrivning av att ``integrera {\"o}ver den del
  av rummet som r{\"o}r v{\aa}rt specifika problem och i den approximation
  som vi identifierat''.}
$$
  q_{\rm tot}
    =\iiint_{{\Bbb R}^3}\rho({\bf x})\,dV
    =q\underbrace{
         \iiint_{{\Bbb R}^3}
           \underbrace{
             \delta({\bf x}-{\bf x}')
           }_{=0,\ {\bf x}\ne{\bf x}'}
           \,dV
       }_{=1}
    =q
$$
Delta-distributionen g{\"o}r h{\"a}r exakt det vi {\"o}nskar, n{\"a}mligen att
{\it lokalisera laddningen till en punkt}, medan faktorn $q$ best{\"a}mmer
v{\aa}r laddnings styrka.

Denna modell {\"a}r sj{\"a}lvfallet i grund och botten endast en bekv{\"a}m
men synnerligen anv{\"a}ndbar idealisering. En verklig elektrisk laddning kan
vara associerad med ett objekt som har en {\"a}ndlig utstr{\"a}ckning, men om
objektet {\"a}r litet j{\"a}mf{\"o}rt med de avst{\aa}nd vi studerar i v{\aa}rt
specifika problem, s{\aa} kan punktladdningsmodellen vara mycket anv{\"a}ndbar.

\subsection{Hur kommer delta-distributionen in i Gauss lag?}
Vi har tidigare sett att Gauss lag f{\"o}r elektrostatiska f{\"a}lt i
differentialform, i omr{\aa}det f{\"o}r en laddningst{\"a}thet
$\rho=\rho({\bf x})$, {\"a}r
$$
  \nabla\cdot{\bf E}=\rho/\varepsilon_0,
$$
och om laddningen i problemet ist{\"a}llet f{\"o}r en kontinuerlig
laddningsf{\"o}rdelning $\rho$ ytg{\"o}rs av en punktladdning $q$ placerad
vid punkten ${\bf x}'$, s{\aa} f{\aa}r vi att
$$
  \nabla\cdot{\bf E}
    =\underbrace{
        q\delta({\bf x}-{\bf x}')
     }_{=\rho({\bf x})}/\varepsilon_0.
$$
Denna synnerligen enkla relation i elektrostatik s{\"a}ger oss direkt ett par
intressanta saker:
\quote{%
  \item{1.}{Vid punkten ${\bf x}'$ finns en lokaliserad {\it k{\"a}lla}
    med total styrka $q/\varepsilon_0$.}
  \item{2.}{{\"O}verallt utom vid k{\"a}llpunkten ${\bf x}'$ {\"a}r
    divergensen f{\"o}r det elektrostatiska f{\"a}ltet noll.}
  \item{3.}{Om vi integrerar divergensen $\nabla\cdot{\bf E}$ {\"o}ver en
    volym som inneh{\aa}ller punkten ${\bf x}'$ erh{\aa}ller vi
    $q/\varepsilon_0$.}
}
\noindent
Med hj{\"a}lp av Gauss teorem (divergensteoremet) f{\"o}ljer d{\"a}rmed direkt
att f{\"o}r varje sluten yta $S$ som omsluter punktladdningen $q$ och som inte
omsluter n{\aa}gon annan laddning, s{\aa} g{\"a}ller det att
$$
  \oiint_S{\bf E}\cdot d{\bf S} = {{q}\over{\varepsilon_0}}.
$$
Det {\"a}r allts{\aa} m{\"o}jligt f{\"o}r ett f{\"a}lt att ha noll divergens i
varje punkt utanf{\"o}r k{\"a}llpunkten, samtidigt som det har ett {\"a}ndligt
totalt fl{\"o}de genom en sluten yta som omsluter k{\"a}llan.
Diracs delta-distribution beskriver rent matematiskt hur detta g{\aa}r ihop
med hj{\"a}lp av konceptet f{\"o}r distributioner.

\subsection{Laddningsf{\"o}rdelningar i en, tv{\aa} och tre dimensioner}
Punktladdningen {\"a}r naturligtvis bara ett av flera exempel p{\aa} hur
delta-funktionen kan anv{\"a}ndas f{\"o}r att beskriva laddningar som {\"a}r
lokaliserade p{\aa} olika s{\"a}tt.
\medskip
\centerline{%
  \epsfbox{figs/dirac-00.1}%
  \hskip 30pt
  \epsfbox{figs/dirac-01.1}%
  \hskip 30pt
  \epsfbox{figs/dirac-02.1}%
}
\medskip
\noindent
Vi kan i huvudsak skilja mellan tre idealiserade fall:
\quote{%
  \item{1.}{{\it Punktladdning}, f{\"o}r vilken laddningen $q$ {\"a}r
    koncentrerad till k{\"a}llpunkten ${\bf x}'$ i tre dimensioner, med
    $$
      \rho({\bf x})=q\delta({\bf x}-{\bf x}')
        =q\delta(x-x')\delta(y-y')\delta(z-z').
    $$
    Den totala laddningen {\"a}r $q$.}
\item{2.}{{\it Linjeladdning}, f{\"o}r vilken laddningen {\"a}r lokaliserad
  till en linje eller kurva. F{\"o}r en rak linje, exempelvis l{\"a}ngs
  $z$-axeln, med linjeladdningst{\"a}thet $\lambda$ (${\rm C}/{\rm m}$),
  har vi laddningst{\"a}theten $\rho({\bf x})=\lambda\delta(x-x')\delta(y-y')$.
  Den totala laddningen {\"a}r $\lambda\times\hbox{linjens l{\"a}ngd}$.
  Integrerar vi {\"o}ver en tv{\"a}rsnittsyta som sk{\"a}r linjen s{\aa}
  erh{\aa}ller vi laddningen per l{\"a}ngdenhet.}
\item{3.}{{\it Ytladdning}, f{\"o}r vilken laddningen {\"a}r koncentrerad till
  en yta. F{\"o}r ett plan, exempelvis $z=0$, med ytladdningst{\"a}theten
  $\sigma$ (${\rm C}/{\rm m}^2$) har vi laddningst{\"a}theten
  $\rho({\bf x})=\sigma\delta(z-z')$. Den totala laddningen {\"a}r
  $\sigma\times\hbox{(planets area)}$. Integrerar vi {\"o}ver en tunn volym
  som omfattar en del av planet f{\aa}r vi laddningen p{\aa} den del av planet
  som volymen omfattar.}
}
\vfill\eject
\noindent
L{\"a}gg m{\"a}rke till m{\"o}nstret f{\"o}r hur laddningst{\"a}theten
$\rho({\bf x})$ med hj{\"a}lp av delta-distributionen kan konstrueras
f{\"o}r punkt-, linje-, och yt-laddningar, med
\medskip
\item{$\bullet$}{Punktladdning: tre delta-distributioner som faktorer, en
  f{\"o}r varje kartesisk koordinat.}
\item{$\bullet$}{Linjeladdning: tv{\aa} delta-distributioner som faktorer,
  eftersom laddningen {\"a}r lokaliserad i tv{\aa} tv{\"a}rg{\aa}ende
  riktningar men kan str{\"a}cka sig l{\"a}ngs den tredje.}
\item{$\bullet$}{Ytladdning: en delta-distribution som faktor, eftersom
  laddningen {\"a}r lokaliserad till ett plan som str{\"a}cker sig i tv{\aa}
  oberoende riktningar.}
\medskip
\noindent
Detta {\"a}r ett mycket anv{\"a}ndbart s{\"a}tt att t{\"a}nka. Antalet
delta-distributioner som ing{\aa}r som faktorer talar om hur m{\aa}nga
riktningar laddningen {\"a}r koncentrerad i. Observera att linjeladdningen
$\lambda$ och ytladdningen $\sigma$ sj{\"a}lvallet kan {\it variera} l{\"a}ngs
linjen respektive {\"o}ver ytan. Uttrycken ovan beskriver endast geometrin
hos laddningsf{\"o}rdelningen i sig, inte n{\"o}dv{\"a}ndigtvis en konstant
laddningst{\"a}thet.

\subsection{Str{\"o}mmar beskrivna med delta-distributioner}
Hittills har vi anv{\"a}nt delta-distributionen f{\"o}r att beskriva elektriska
laddningar, dock kan den lika v{\"a}l anv{\"a}ndas f{\"o}r att beskriva
elektriska str{\"o}mmar.
Str{\"o}mt{\"a}theten ${\bf J}$ anger str{\"o}m per tv{\"a}rsnittsarea och har
enheten ${\rm A}/{\rm m}^2$, samt f{\"o}rekommer i bland annat Amp{\`e}res
magnetostatiska lag\sidx{Amp{\`e}res lag}[Magnetostatik]
$$
  \nabla\times{\bf B}=\mu_0{\bf J}
$$
och, om vi l{\"a}gger till f{\"o}rskjutningsstr{\"o}mmen som vi kommer att
behandla i detalj i F{\"o}rel{\"a}sning~9, i Amp{\`e}re--Maxwells
elektrodynamiska lag
$$
  \nabla\times{\bf B}=\mu_0{\bf J}
    +\mu_0\varepsilon_0{{\partial{\bf E}}\over{\partial t}}
$$\sidx{Amp{\`e}res lag}[Amp{\`e}re--Maxwells elektrodynamiska lag]
Antag nu att en tunn ledare l{\"o}per l{\"a}ngs $z$-axeln och att str{\"o}mmen
$I$ flyter i positiv $z$-riktning i ledaren. Om den tunna ledaren {\"a}r
placerad i koordinaten $(x',y')$, bara f{\"o}r att h{\aa}lla oss till v{\aa}r
konvention med primmade k{\"a}llkoordinater, s{\aa} kan vi beskriva denna
ledare som b{\"a}randes en str{\"o}mt{\"a}thet beskriven i delta-distributioner
$$
  {\bf J}({\bf x}) =I\delta(x-x')\delta(y-y'){\bf e}_z.
$$
Str{\"o}mmen {\"a}r i detta uttryck koncentrerad till punkten $(x,y)=(x',y')$
i tv{\"a}rsnittet $z=\hbox{konstant}$, men str{\"a}cker sig l{\"a}ngs hela
$z$-axeln.\sidx{Str{\"o}mt{\"a}thet ${\bf J}$}
I detta fall beh{\"o}ver vi med andra ord tv{\aa} delta-distributioner f{\"o}r
att s{\aa} att s{\"a}ga ``h{\aa}lla ihop'' str{\"o}mmen, f{\"o}r de tv{\aa}
tv{\"a}rg{\aa}ende koordinaterna $x$ och $y$.
\sidx{Str{\"o}mslinga}
\sidx{Str{\"o}mslinga}[Delta-distribution och str{\"o}mt{\"a}thet]

F{\"o}r att kontrollera uttrycket f{\"o}r str{\"o}mmen l{\"a}ngs en linje
enligt ovan, s{\aa} kan vi integrera str{\"o}mt{\"a}theten {\"o}ver en
tv{\"a}rsnittsyta som {\"a}r vinkelr{\"a}t mot $z$-axeln, enligt
$$
  \eqalign{
    I_{\rm tot}=\iint_{{\Bbb R}^2}{\bf J}({\bf x})\cdot{\bf e}_z\,dy\,dx
      =I\underbrace{\int\delta(x-x')\,dx}_{=1}
        \underbrace{\int\delta(y-y')\,dy}_{=1}
      =I,
  }
$$
med andra ord f{\aa}r vi allts{\aa} tillbaka den totala str{\"o}mmen $I$,
precis som v{\"a}ntat. Detta f{\"o}ljer exakt samma princip som tidigare
f{\"o}r punktladdningen, men att vi nu st{\"a}llet beskriver en str{\"o}m
som {\"a}r koncentrerad till en linje.

\subsection{En pulsad str{\"o}m}
Vi kan ocks{\aa} beskriva en str{\"o}m som {\"a}r koncentrerad till en viss
punkt i tiden, i form av en idealiserad impuls. Om vi som vanligt l{\aa}ter
$t$ beteckna tiden, och l{\aa}ter str{\"o}mmen $I(t)$ beskrivas av en puls
vid tiden $t=t'$ som\sidx{Str{\"o}m}[Pulsad]
\sidx{Str{\"o}m}[Transporterad laddning]
$$
  I(t)=q\delta(t-t'),
$$
s{\aa} g{\"a}ller, i analogi med tidigare problemformulering i spatial
dom{\"a}n, att den temporala\numberedfootnote{$\hbox{``Spatialt''}
  =\hbox{i rummet}$; $\hbox{``Temporalt''}=\hbox{i tiden}$.}
pulsen transporterar den totala laddningen
$$
  \int^{\infty}_{-\infty}I(t)\,dt=q\int^{\infty}_{-\infty}\delta(t-t')\,dt=q,
$$
d{\"a}r $I(t)$ har enheten ampere och $\delta(t-t')$ ist{\"a}llet enheten
$1/{\rm s}$. Detta {\"a}r en idealiserad str{\"o}mpuls med o{\"a}ndligt kort
varaktighet men {\"a}ndlig {\"o}verf{\"o}rd laddning. En verklig str{\"o}mpuls
har naturligtvis en {\"a}ndlig varaktighet och en begr{\"a}nsad toppstr{\"o}m.

\subsection{Kontinuitetsekvationen f{\"o}r elektrisk laddning}
Det finns en viktig fysikalisk begr{\"a}nsning n{\"a}r vi v{\"a}ljer laddnings-
och str{\"o}mf{\"o}rdelningar, n{\"a}mligen att {\it elektrisk laddning alltid
bevaras.} \sidx{Kontinuitetsekvationen} Detta centrala begrepp ligger bland
annat bakom den f{\"o}rskjut\-nings\-str{\"o}m\numberedfootnote{Mer om
  f{\"o}rskjutningsstr{\"o}mmen och dess ursprung som term i
  Amp\`ere--Maxwells ekvation i F{\"o}rel{\"a}sning~9, d{\aa}
  vi g{\aa}r igenom Maxwells ekvationer.}
som formulerades av James Clerk Maxwell \sidx{Maxwell, James Clerk (1831--1879)}
som bas f{\"o}r den teori som beskriver elektromagnetisk v{\aa}gutbredning, och
som uttrycks av kontinuitetsekvationen f{\"o}r laddning,
$$
  {{\partial\rho}\over{\partial t}}+\nabla\cdot{\bf J}=0.
$$
Kontinuitetsekvationen s{\"a}ger oss att {\it nettoladdningen i ett omr{\aa}de
endast kan minska om det finns ett nettofl{\"o}de av laddning ut genom
omr{\aa}dets rand, och omv{\"a}nt.}
Detta inneb{\"a}r i sin tur att vi inte kan v{\"a}lja laddningst{\"a}theten
$\rho$ och str{\"o}mt{\"a}theten ${\bf J}$ helt oberoende av varandra.
Om laddningst{\"a}theten f{\"o}r{\"a}ndras med tiden s{\aa} m{\aa}ste
str{\"o}mmen vara f{\"o}renlig med denna f{\"o}r{\"a}ndring.
Ett brott mot kontinuitetsekvationen {\"a}r ocks{\aa} ett brott mot (den
klassiska) fysiken, och skulle i s{\aa} fall inneb{\"a}ra att laddning
pl{\"o}tsligt skapas eller f{\"o}rintas.

Exempelvis kan en punktladdning som r{\"o}r sig beskrivas med
laddningst{\"a}theten\numberedfootnote{I detta fall har vi att g{\"o}ra
  med en delta-distribution av ett argument i tre dimensioner (en vektor
  i rummet), som d{\"a}rmed har enheten $1/{\rm m}^3$; med andra ord har
  produkten $q\delta^{(3)}$ enheten ${\rm C}/{\rm m}^3$, precis som v{\"a}ntat.}
$$
  \rho({\bf x},t)=q\delta^{(3)}({\bf x}-{\bf x}_q(t))
$$
d{\"a}r ${\bf x}_q(t))$ {\"a}r laddningens r{\"o}rliga position. Den motsvarande
str{\"o}mt{\"a}theten blir i detta fall
$$
  {\bf J}({\bf x},t)=q{\bf v}_q(t) \delta^{(3)}\big({\bf x}-{\bf x}_q(t)\big)
$$
d{\"a}r ${\bf v}_q(t)=d{\bf x}_q/dt$ {\"a}r laddningens hastighet. Som vi ser
i detta exempel, s{\aa} {\"a}r laddningst{\"a}theten och str{\"o}mt{\"a}theten
koncentrerade till samma r{\"o}rliga punkt. Dessa uttryck uppfyller
kontinuitetsekvationen i distributionell mening, och utg{\"o}r ett exempel
p{\aa} hur delta-distributionen l{\aa}ter oss beskriva en idealiserad fysisk
k{\"a}lla p{\aa} ett matematiskt konsekvent s{\"a}tt.\numberedfootnote{Det
  finns dock ett grundl{\"a}ggande problem med detta exempel, som g{\"o}r
  att vi m{\aa}ste vara f{\"o}rsiktiga i att dra n{\aa}gra djupare
  slutsatser fr{\aa}n resonemang kring transporten av en enda punktladdning,
  n{\"a}mligen att detta inte p{\aa} n{\aa}got s{\"a}tt utg{\"o}r en
  kontinuerlig str{\"o}m i den bem{\"a}rkelse som vi normalt anv{\"a}nder
  inom elektromagnetisk teori. En punktladdning som r{\"o}r sig {\"a}r
  till sin natur s{\aa}dan att den ena {\"o}gonblicket befinner sig p{\aa}
  en position, och {\"o}gonblicket d{\"a}refter {\"a}r borta.
  Detta {\"a}r inte en grund som exempelvis Amp{\`e}res eller Biot-Savarts
  lagar kan formuleras utifr{\aa}n, och det {\"a}r ocks{\aa} en
  inv{\"a}ndning som Griffiths formulerar i {\it Introduction to
  Electrodynamics}. Mer om detta i avsnittet om Biot--Savarts lag
  i F{\"o}rel{\"a}sning~5.}
\vfill\eject

\subsection{Att ta med sig fr{\aa}n delta-distributionen}
Innan vi g{\aa}r vidare till vektoranalysen och begrepp som gradient, divergens
och rotation, s{\aa} kan vi passa p{\aa} att sammanfatta tre grundl{\"a}ggande
begrepp kring delta-distributionen.
\quote{%
  \item{1.}{Diracs delta-distribution beskriver i rummet en {\it idealiserad
    lokal k{\"a}lla}. Denna k{\"a}lla {\"a}r koncentrerad till en punkt, en
    linje, en yta eller en tidpunkt, beroende p{\aa} hur den definieras
    utifr{\aa}n kombinationer av delta-distributioner.
    \sidx{Diracs delta-distribution $\delta({\bf x})$}}
  \item{2.}{Delta-distributionens {\it styrka best{\"a}ms av faktorn
    framf{\"o}r den}. F{\"o}r en punktladdning {\"a}r denna faktor $q$
    (laddningen i sig), medan faktorn f{\"o}r en linjeladdning {\"a}r
    $\lambda$ (laddning per l{\"a}ngdenhet) och f{\"o}r en ytladdning
    $\sigma$ (laddning per ytenhet). Genom att integrera {\"o}ver
    l{\"a}mplig dom{\"a}n f{\aa}r vi tillbaka den totalt inneslutna
    laddningen eller str{\"o}mmen genom en yta.}
  \item{3.}{Delta-distributionen {\"a}r {\it s{\"a}rskilt anv{\"a}ndbar
    n{\"a}r vi vill beskriva f{\"a}lt med lokaliserade k{\"a}llor}.
    Gauss lag kan exempelvis skrivas som
    $$
      \nabla\cdot{\bf E}
        =\underbrace{
            q\,\delta^{(3)}({\bf x}-{\bf x}')
         }_{=\rho({\bf x})}/\varepsilon_0
    $$
    {\"a}ven f{\"o}r en punktladdning placerad i k{\"a}llpunkten ${\bf x}'$.
    H{\"a}r blir kopplingen mellan divergensen $\nabla\cdot{\bf E}$ och
    k{\"a}llan $q\,\delta^{(3)}({\bf x})/\varepsilon_0$ uppenbar.}
}
\noindent
Delta-distributionen $\delta({\bf x})$ kan allts{\aa} ses som {\it bryggan
mellan en idealiserad k{\"a}lla och de differentialekvationer som beskriver
dess f{\"a}lt}.

N{\"a}r vi nu g{\aa}r vidare till vektoranalysen kan vi fr{\aa}ga oss hur
f{\"a}ltet ser ut i omgivningen av en s{\aa}dan k{\"a}lla, var divergensen
{\"a}r noll, och var eventuella singulariteter finns? Vidare, hur kan vi
beskriva f{\"a}ltets riktning och variation med gradient och rotation, och
hur kommer vi fram till kopplingar mellan f{\"a}lt?
Kort och gott, det {\"a}r nu den egentliga vektoranalysen och byggandet av
v{\aa}r ``matematiska verktygsl{\aa}da'' b{\"o}rjar.
\vfill\eject

\section{{\hl{Gradienten $\nabla\phi$}}}
\subsection{Tolkning av gradienten i ett landskap best{\aa}ende av
  kullar och dalar}
F{\"o}rest{\"a}ll dig att du st{\aa}r i ett kuperat landskap, d{\"a}r
h{\"o}jden {\"o}ver havet varierar fr{\aa}n punkt till punkt.
Vi kan beskriva h{\"o}jden med ett skal{\"a}rf{\"a}lt $h=h({\bf x})$,
d{\"a}r vi f{\"o}r enkelhets skull kan se ${\bf x}=(x,y)$ som de sedvanliga
tv{\aa}-dimensionella kartkoordinaterna i l{\"a}mpligt valt koordinatsystem.
\epsfig{figs/potmap.1}
\noindent
Om vi {\"o}nskar ta reda p{\aa} hur h{\"o}jden f{\"o}r{\"a}ndras allt eftersom
vi r{\"o}r oss i landskapet i en viss riktning, s{\aa} kan vi ber{\"a}kna en
riktningsderivata, exempelvis som $\partial/\partial x$ om vi {\"o}nskar finna
f{\"o}r{\"a}ndringen l{\"a}ngs den $x$-axel som vi definierat enligt v{\aa}r
karta.
Om vi ist{\"a}llet {\"o}nskar ta reda p{\aa} i vilken riktning h{\"o}jden
{\it {\"o}kar allra snabbast}, och {\it hur snabbt} den {\"o}kar, s{\aa}
kan vi anv{\"a}nda gradienten av denna ``potential'' $h({\bf x})$, som
$$
  \nabla h={{\partial h}\over{\partial x}}{\bf e}_x
    +{{\partial h}\over{\partial y}}{\bf e}_y,
$$
som allts{\aa} ger oss en vektor som pekar i riktningen f{\"o}r den snabbaste
{\it {\"o}kningen} av h{\"o}jden, och d{\"a}r dess belopp anger den maximala
h{\"o}jd{\"o}kningen per l{\"a}ngdenhet.
Om vi st{\aa}r p{\aa} en sluttning pekar gradienten uppf{\"o}r sluttningen
mot den brantaste stigningen, och omv{\"a}nt g{\"a}ller att samma vektor med
omv{\"a}nt tecken, $-\nabla h$, pekar {\it nedf{\"o}r} den brantaste
sluttningen.

Vi kan ocks{\aa} t{\"a}nka oss h{\"o}jdkurvor p{\aa} en karta. L{\"a}ngs en
h{\"o}jdkurva {\"a}r h{\"o}jden konstant, medan vektorn f{\"o}r gradienten
alltid {\"a}r riktad ortogonalt mot h{\"o}jdkurvan eftersom den pekar i den
riktning d{\"a}r h{\"o}jden {\"a}ndras snabbast.
Detta bist{\aa}r oss med en mycket anv{\"a}ndbar geometrisk tolkning,
n{\"a}mligen att
\quote{%
  {\it gradienten {\"a}r riktad vinkelr{\"a}tt mot ytor eller kurvor
  d{\"a}r skal{\"a}rf{\"a}ltet {\"a}r konstant.}
}

\subsection{Gradientens matematiska definition}
Bilden av gradienten i tv{\aa} dimensioner som den vektor som lokalt pekar ut
den brantaste sluttningen {\aa}t oss {\"a}r en bra sinnebild som vi kan ta med
oss. I fallet med elektrodynamik s{\aa} har vi dock alltsom oftast att g{\"o}ra
med problem i tre dimensioner, och vi beh{\"o}ver i och med detta en g{\aa}ng
f{\"o}r alla formalisera gradienten.

I kartesiska koordinater definierar vi gradienten i tre dimensioner av ett
skal{\"a}rf{\"a}lt $\phi({\bf x})$ som
$$
  \hl{
  \nabla\phi
    \equiv{\bf e}_x{{\partial\phi}\over{\partial x}}
        +{\bf e}_y{{\partial\phi}\over{\partial y}}
        +{\bf e}_z{{\partial\phi}\over{\partial z}}
        }
    =\Big(
       {{\partial\phi}\over{\partial x}},
       {{\partial\phi}\over{\partial y}},
       {{\partial\phi}\over{\partial z}}
     \Big)
$$
Operatorn $\nabla$, som uttalas ``nabla'' eller ``dell'', {\"a}r kort och gott
ingenting annat {\"a}n en vektor av partiella derivator. N{\"a}r vi applicerar
nabla-operatorn $\nabla$ p{\aa} ett skal{\"a}rt f{\"a}lt $\phi$ erh{\aa}ller
vi en vektor best{\aa}ende av gradienten ${\rm grad}\phi$ av f{\"a}ltet.

\subsection{Exempel: elektrisk potential och elektriskt f{\"a}lt}
L{\aa}t oss nu byta ut landskapets h{\"o}jd mot den elektriska potential i
tre dimensioner, ekvivalent med den skal{\"a}ra potentialen $\phi({\bf x})$.
Antag att potentialen mellan, s{\"a}g, tv{\aa} elektroder varierar med
positionen mellan dem. Om vi k{\"a}nner hur den skal{\"a}ra potentialen
$\phi({\bf x})$ varierar i rummet, s{\aa} inneb{\"a}r det att vi direkt kan
best{\"a}mma det elektrostatiska f{\"a}ltet genom
$$
  {\bf E}=-\nabla\phi.
$$
I ord betyder detta att det {\it elektriska f{\"a}ltet pekar i riktningen
f{\"o}r den snabbaste minskningen av den skal{\"a}ra potentialen}.
\epsfig{figs/potcapac.1}
\noindent
F{\"o}r att ta ett enkelt exempel, betrakta en skal{\"a}r potential som endast
varierar l{\"a}ngs en enda riktning, vilken vi h{\"a}r kan definiera som
$x$-axeln, med
$$
  \phi(x)=\phi_0-({{U_0}/{a}})x
$$
d{\"a}r $U_0$ och $a$ {\"a}r positiva konstanter som betecknar skillnad i
elektrisk potential (${\rm V}$) och elektrodavst{\aa}nd (${\rm m}$).
Ett typiskt exempel f{\"o}r denna linj{\"a}ra variation i potential {\"a}r
omr{\aa}det mellan plattorna i en vanlig plattkondensator, d{\"a}r vi ignorerar
randeffekter.
I detta exempel erh{\aa}ller vi direkt att det elektriska f{\"a}ltet {\"a}r
$$
  {\bf E}
    =-\nabla\phi
    =-{\bf e}_x{{\partial}\over{\partial x}}
        \underbrace{\big(\phi_0-({{U_0}/{a}})x\big)}_{=\phi(x)}
    ={\bf e}_xU_0/a,
$$
med andra ord har vi allts{\aa} att g{\"o}ra med ett {\it homogent} elektriskt
f{\"a}lt som pekar i positiv $x$-riktning, och som v{\"a}ntat f{\"o}r ett
f{\"a}lt mellan tv{\aa} plattor med potentialskillnad $U_0$, separerade med
ett avst{\aa}nd $a$, s{\aa} blir ocks{\aa} beloppet av det elektriska
f{\"a}ltet just $E=U_0/a$.
Den skal{\"a}ra potentialen minskar i f{\"a}ltets riktning, helt i analogi med
v{\aa}rt tidigare exempel med h{\"o}jdkurvor i ett landskap, d{\"a}r h{\"o}jden
(potentialen) minskar n{\"a}r vi g{\aa}r nedf{\"o}r den brantaste sluttningen
om vi f{\"o}ljer $-\nabla h$.
\quote{%
  {\it Notera att f{\"o}r en konstant skal{\"a}r potential $\phi$ {\"a}r
  det elektriska f{\"a}ltet alltid identiskt noll!}\numberedfootnote{Detta
    {\"a}r grunden f{\"o}r den s{\aa} kallade ``Faradays bur'', som vi
    kommer att g{\aa} igenom i detalj i F{\"o}rel{\"a}sning~3.}
  {\it Utan en variation i skal{\"a}r potential kan vi inte ha ett nollskilt
  elektrostatiskt f{\"a}lt.}
}
\noindent
F{\"o}r en positiv testladdning $q$ minskar den potentiella energin n{\"a}r
laddningen r{\"o}r sig i det elektriska f{\"a}ltets riktning. Eftersom den
elektriska potentiella energin {\"a}r $U=q\phi$, s{\aa} g{\"a}ller det att
en f{\"o}r{\"a}ndring av den skal{\"a}ra potentialen direkt proportionellt
motsvaras av en f{\"o}r{\"a}ndring av den elektrostatiskt upplagrade energin
hos laddningen,
$$
  \Delta U=q\,\Delta \phi,
$$
det vill s{\"a}ga att f{\"o}r en positiv laddning $q$ s{\aa} inneb{\"a}r en
minskning $\Delta\phi$ av den skal{\"a}ra potentialen ocks{\aa} en minskning
$\Delta U$ av den potentiella energin.

Rent ``f{\"a}ltm{\"a}ssigt'' {\"a}r det h{\"a}r viktigt att skilja mellan
potentialen $\phi$ i sig och det associerade elektrostatiska f{\"a}ltet
${\bf E}$, d{\"a}r potentialen {\"a}r ett {\it skal{\"a}rf{\"a}lt}, medan
det elektriska f{\"a}ltet {\"a}r ett {\it vektorf{\"a}lt} som beskriver
kraften per laddningsenhet f{\"o}r en hypotetisk testladdning.
Notera att det {\"a}ven {\"a}r mycket l{\"a}tt att l{\"a}ngs med resans
g{\aa}ng f{\"o}rv{\"a}xla den skal{\"a}ra potentialen $\phi$, som m{\"a}ts
i Volt, med potentiell energi, som m{\"a}ts i Joule.

\section{{\hl{Divergensen $\nabla\cdot{\bf a}$}}}
\subsection{Vad beskriver divergensen?}
Vi l{\"a}mnar nu skal{\"a}ra f{\"a}lt och g{\aa}r {\"o}ver till vektorf{\"a}lt.
F{\"o}rest{\"a}ll dig att du st{\aa}r mitt i en str{\"o}mmande v{\"a}tska,
d{\"a}r v{\"a}tskan kan str{\"o}mma in mot en punkt, ut fr{\aa}n en punkt
eller passera f{\"o}rbi utan att det sker n{\aa}gon lokal ansamling eller
uttunning.
\quote{%
  {\it Divergensen {\"a}r ett m{\aa}tt p{\aa} det lokala nettofl{\"o}det
  ut fr{\aa}n en punkt.}
}
\noindent
F{\"o}r ett vektorf{\"a}lt ${\bf a}$ definieras divergensen som
$$
  {\rm div}({\bf a})\equiv\nabla\cdot{\bf a}.
$$
I kartesiska koordinater {\"a}r divergensen allts{\aa} given som
$$
  \eqalign{
  \hl{\nabla\cdot{\bf a}}&\hl{\equiv
    {{\partial a_x}\over{\partial x}}
        +{{\partial a_y}\over{\partial y}}
        +{{\partial a_z}\over{\partial z}}}\cr
    &=\Big(
        {\bf e}_x{{\partial}\over{\partial x}}
          +{\bf e}_y{{\partial}\over{\partial y}}
          +{\bf e}_z{{\partial}\over{\partial z}}
      \Big)\cdot\Big({\bf e}_xa_x+{\bf e}_ya_y+{\bf e}_za_z\Big).\cr
  }
$$
Observera att divergensen {\"a}r ett skal{\"a}rt m{\aa}tt som funktion {\"o}ver
rummet, och att divergensen utg{\"o}r ett skal{\"a}rt f{\"a}lt i sig.
Vi b{\"o}rjar med en vektor ${\bf a}$ och erh{\aa}ller fr{\aa}n divergensen
ett tal i varje punkt, d{\"a}r detta tal kan tolkas enligt f{\"o}ljande, med
analogi till v{\"a}tskefl{\"o}det,
\quote{%
  \item{\hl{$\bullet$}}{\hl{{\it Positiv divergens} inneb{\"a}r att
    det lokalt finns ett {\it nettofl{\"o}de ut{\aa}t}. Vektorf{\"a}ltet
    ${\bf a}$ beter sig som att det i punkten finns en {\it k{\"a}lla}.}}
  \item{\hl{$\bullet$}}{\hl{{\it Negativ divergens} inneb{\"a}r att
    det lokalt finns ett {\it nettofl{\"o}de in{\aa}t}. Vektorf{\"a}ltet
    ${\bf a}$ beter sig som att det i punkten finns en {\it s{\"a}nka}.}}
  \item{\hl{$\bullet$}}{\hl{Noll divergens inneb{\"a}r att nettofl{\"o}det
    in{\aa}t och ut{\aa}t fr{\aa}n punkten {\"a}r noll.
    Detta betyder {\it inte} n{\"o}dv{\"a}ndigtvis att inget fl{\"o}de
    finns, utan att infl{\"o}de och utfl{\"o}de balanserar varandra.}%
    \numberedfootnote{Exempelvis om vi t{\"a}nker oss att vi
      st{\aa}r i en v{\"a}tskestr{\"o}m i en {\aa}, d{\"a}r vattnet
      str{\"o}mmar f{\"o}rbi oss utan att det finns n{\aa}gra lokala
      vattenk{\"a}llor eller s{\"a}nkor. Divergensen i punkten vi
      st{\aa}r i {\"a}r noll, trots ett h{\"o}gst n{\"a}rvarande fl{\"o}de.}}
}
\noindent
\epsfig{figs/field-divzero.1}
\noindent
Ett enkelt exempel p{\aa} ett vektorf{\"a}lt med noll divergens {\"a}r ett
homogent vektorf{\"a}lt, med
$$
  {\bf a}=a_0{\bf e}_x,
$$
f{\"o}r vilket komponenterna {\"a}r konstanta och d{\"a}rmed har att alla
derivator {\"a}r identiskt noll. D{\"a}rf{\"o}r g{\"a}ller f{\"o}r ett homogent
vektorf{\"a}lt att
$$
  \nabla\cdot{\bf a}=0.
$$
Pilarna i det homogena vektorf{\"a}ltet {\"a}r lika l{\aa}nga och lika riktade
{\"o}verallt. Inget lokalt omr{\aa}de har ett nettofl{\"o}de ut{\aa}t eller
in{\aa}t.
\epsfig{figs/field-div-nonzero.1}
\noindent
J{\"a}mf{\"o}r detta homogena exempel med det {\it inhomogena} vektorf{\"a}ltet
$$
  {\bf a}=x{\bf e}_x+y{\bf e}_y+z{\bf e}_z,
$$
f{\"o}r vilket vi ist{\"a}llet erh{\aa}ller att
$$
  \nabla\cdot{\bf a}=1+1+1=3.
$$
F{\"a}ltet pekar i detta inhomogena fall radiellt ut{\aa}t fr{\aa}n origo och
v{\"a}xer dessutom med avst{\aa}ndet. {\it Divergensen {\"a}r i detta exempel
positiv {\"o}verallt, och notera speciellt att detta {\"a}ven g{\"a}ller i
origo, trots att f{\"a}ltet d{\"a}r {\"a}r noll!}

\subsection{Divergens och fl{\"o}de genom en sluten yta}
Vi kan f{\"o}rst{\aa} konceptet med divergensen {\"a}nnu b{\"a}ttre genom att
betrakta ett litet omr{\aa}de som omges av en sluten yta, ut genom vilken vi
nu kommer att studera det s{\aa} kallade {\it fl{\"o}det}.
{\hl{Fl{\"o}det av ett vektorf{\"a}lt}} ut genom en sluten yta $S$ definieras
som\numberedfootnote{Notera den subtila skillnaden i notation h{\"a}r,
  d{\"a}r vi genomg{\aa}ende i f{\"o}rel{\"a}sningsserien l{\aa}ter
  ``$\Phi$'' beteckna ett {\it fl{\"o}de}, medan ``$\phi$'' betecknar
  den {\it skal{\"a}ra potentialen}. F{\"o}r att ytterligare riskera
  att {\"o}ka f{\"o}rvirringen kommer vi dessutom att senare l{\aa}ta
  ``$\varphi$'' beteckna fasen hos en v{\aa}g.}
$$
  \hl{\Phi=\oiint{\bf a}_S\cdot d{\bf S}},
$$
d{\"a}r $d{\bf S}$ {\"a}r ett ut{\aa}triktat ytelement, med en riktning som
lokalt {\"a}r normal mot ytan $S$.

Gauss lag, {\"a}ven kallad {\it divergensteoremet}, s{\"a}ger att det totala
ut{\aa}triktade fl{\"o}det genom en sluten yta {\"a}r lika med volymintegralen
av divergensen innanf{\"o}r ytan, eller
$$
  \oiint_S{\bf a}\cdot d{\bf S}=\iiint_V(\nabla\cdot{\bf a})\,dV.
$$
Vi kommer i denna kurs att mycket flitigt anv{\"a}nda oss av Gauss lag, och
detta {\"a}r en grundbult inom elektromagnetisk f{\"a}ltteori som {\"a}r
v{\"a}l v{\"a}rd att l{\"a}ra sig utantill.

Ut{\"o}ver att Gauss lag bist{\aa}r oss med assistans i rent praktiskt
probleml{\"o}sande, s{\aa} {\"a}r det dessutom utifr{\aa}n denna lag som
kopplingen mellan divergens och s{\aa} kallade {\it k{\"a}llor}, eller
{\it k{\"a}lltermer}, inom elektromagnetism kommer, i analogi med det tidigare
tanke-experimentet kring v{\"a}tskefl{\"o}de och divergens.
{\it Om den totala divergensen i ett omr{\aa}de {\"a}r positivt m{\aa}ste det
finnas ett nettofl{\"o}de ut genom omr{\aa}dets rand}, och det {\"a}r genom
ytintegralen ovan som vi tolkar ett s{\aa}dant ``fl{\"o}de $\Phi$ av
f{\"a}ltlijer'' av, exempelvis, eletriska och magnetiska f{\"a}lt.

Det skall h{\"a}r dock po{\"a}ngteras att divergensen inte
n{\"o}dv{\"a}ndigtvis beskiver en fysisk k{\"a}lla i form av ett
f{\"o}rem{\aa}l. Divergensen beskriver i grund och botten endast en lokal
matematisk egenskap hos f{\"a}ltet i sig. I elektromagnetism kopplar
v{\aa}ra ekvationer denna egenskap till elektrisk laddning.

\subsection{Divergensen hos det elektriska f{\"a}ltet}
En av ekvationerna inom elektromagnetism {\"a}r Gauss lag f{\"o}r det
elektriska f{\"a}ltet,\numberedfootnote{Gauss lag {\"a}r i denna form
  giltig f{\"o}r statiska (tidsoberoende) s{\aa}v{\"a}l som dynamiska
  (tidsvarie\-rande) elektriska f{\"a}lt.}
som p{\aa} differentialform lyder
$$
  \nabla\cdot{\bf E}=\rho/\varepsilon_0,
$$
d{\"a}r $\rho$ {\"a}r den elektriska laddningst{\"a}theten och $\varepsilon_0$
den elektriska permittiviteten i vakuum, eller kort och gott
{\it vakuumpermittiviteten}.
Denna ekvation s{\"a}ger oss att det elektriska f{\"a}ltets divergens
best{\"a}ms av laddningst{\"a}theten, som i sin tur ger oss tre principiella
fall:
\quote{%
  \item{1.}{Om $\rho>0$, s{\aa} {\"a}r divergensen $\nabla\cdot{\bf E}$
    f{\"o}r det elektriska f{\"a}ltet positiv. Vi kan se det som att
    {\it en positiv laddningst{\"a}thet} $\rho$ skapar ett {\it positivt
    fl{\"o}de}, riktat ut{\aa}t, som i sin tur kan tolkas som att en
    positiv laddning svarar mot en {\it k{\"a}lla}.}
  \item{2.}{Om ist{\"a}llet $\rho<0$, s{\aa} {\"a}r divergensen negativ.
    I detta fall kan vi ist{\"a}llet se det som att {it en negativ
    laddningst{\"a}thet} $\rho$ skapar ett {\it negativt fl{\"o}de},
    riktat in{\aa}t, som i sin tur kan tolkas som att en negativ
    laddning svarar mot en {\it s{\"a}nka}.}
  \item{3.}{Om $\rho=0$, s{\aa} {\"a}r divergensen noll.}
}
\noindent
Det tredje fallet kr{\"a}ver en viktig precisering, n{\"a}mligen att ett
elektriskt f{\"a}lt alltid {\"a}r divergensfritt i ett omr{\aa}de som
saknar elektrisk laddningst{\"a}thet,\numberedfootnote{Notera att
  relationen $\nabla\cdot{\bf E}=\rho/\varepsilon_0$ g{\"a}ller
  {\it punktvis}, det vill s{\"a}ga att vi inte n{\"o}dv{\"a}ndigtvis
  beh{\"o}ver innesluta en laddning via en sluten yta $S$. Vi kan precis
  lika g{\"a}rna l{\aa}ta den omslutande ytan $S$ i Gauss lag p{\aa}
  integralform g{\aa} mot noll (en punkt), och i gr{\"a}nsen evaluera
  Gauss lag p{\aa} differentialformen som h{\"a}r anv{\"a}nds.}
till exempel f{\"a}ltet fr{\aa}n en positiv punktladdning $q$ placerad
i k{\"a}llpunkten ${\bf x}'$, d{\"a}r f{\"a}ltlinjerna alla pekar ut{\aa}t.
\epsfig{figs/field-pointcharge.1}
\noindent
Det {\"a}r h{\"a}r l{\"a}tt att f{\"o}rledas att t{\"a}nka att ``eftersom
f{\"a}ltlinjerna utanf{\"o}r laddningen ju faktiskt {\it divergerar}, s{\aa}
borde vi ju ha en divergens''. Detta resonemang {\"a}r dock felaktigt, och vi
rekapitulerar att {\it divergensen talar om f{\"o}r oss var k{\"a}llor och
s{\"a}nkor finns}.
Detta resonemang kring Gauss lag f{\"o}r elektriska f{\"a}lt g{\"a}ller
b{\aa}de i elektrostatik och i elektrodynamik, s{\aa} l{\"a}nge vi
anv{\"a}nder Maxwells klassiska makroskopiska eller mikroskopiska formulering
med motsvarande laddningst{\"a}thet och f{\"a}lt.

Utanf{\"o}r sj{\"a}lva k{\"a}llpunkten ${\bf x}'$ {\"a}r divergensen noll,
medan punktladdningen ger en singularitet i origo. Matematiskt beskriver vi
enligt tidigare detta med en delta-distribution som
$$
  \nabla\cdot{\bf E} ={{q}\over{\varepsilon_0}}\delta^{(3)}({\bf x}).
$$
Detta illustrerar en viktig skillnad mellan ett f{\"a}lt som {\"a}r
divergensfritt i varje punkt i ett omr{\aa}de och ett f{\"a}lt som har
en singular k{\"a}lla i omr{\aa}det.

\subsection{Divergensen hos det magnetiska f{\"a}ltet}
F{\"o}r magnetf{\"a}ltet\numberedfootnote{{\AA}terigen, om vi skall vara
petiga, den {\it magnetiska fl{\"o}dest{\"a}theten} ${\bf B}$.} ${\bf B}$
har vi en annan av Maxwells ekvationer som vi ofta kommer att {\aa}terkomma
till i denna kurs, n{\"a}mligen att {\it divergensen f{\"o}r det magnetiska
f{\"a}ltet, alltid {\"a}r identiskt noll},
$$
  \nabla\cdot{\bf B}=0.
$$
Denna lag kallas kort och gott f{\"o}r {\it Gauss lag f{\"o}r det magnetiska
f{\"a}ltet}.\numberedfootnote{Vi kommer att rent formellt h{\"a}rleda
  divergensen f{\"o}r magnetiska f{\"a}lt utifr{\aa}n Biot--Savarts
  lag i F{\"o}rel{\"a}sning~4.}

Till skillnad fr{\aa}n det elektriska f{\"a}ltet existerar det i den klassiska
teorin f{\"o}r elektromagnetism inga isolerade magnetiska laddningar, eller
{\it magnetiska monopoler}, och divergensen f{\"o}r magnetf{\"a}ltet {\"a}r
d{\"a}rf{\"o}r alltid identiskt noll {\"o}verallt, {\"a}ven i omr{\aa}den
d{\"a}r det finns elektriska str{\"o}mmar eller magnetiska material.
F{\"a}ltlinjerna f{\"o}r magnetf{\"a}ltet har allts{\aa} inga start- eller
{\"a}ndpunkter i form av isolerade magnetiska k{\"a}llor eller s{\"a}nkor.
F{\"a}ltlinjerna f{\"o}r magnetf{\"a}lt bildar alltid slutna slingor och
l{\"o}per genom rummet utan att b{\"o}rja eller sluta i n{\aa}gon
monopolladdning.
Det betyder {\it inte} att magnetf{\"a}ltet saknar struktur eller att det
m{\aa}ste vara homogent. {\it Ett f{\"a}lt kan vara kraftigt varierande och
{\"a}nd{\aa} ha noll divergens.}

\subsection{Konkret exempel p{\aa} att divergensen f{\"o}r f{\"a}ltet
  utanf{\"o}r en punktladdning {\"a}r noll}
Vi vet att det elektriska f{\"a}ltet fr{\aa}n en punktladdning $q$, vilken vi
f{\"o}r enkelhets skull kan definiera som v{\aa}rt origo, {\"a}r ett radiellt
ut{\aa}triktat f{\"a}lt som har ett belopp som g{\aa}r som $1/|{\bf x}|^2$.
Det kan intuitivt k{\"a}nnas lite konstigt att tro att en ``derivata'' av en
s{\aa}dan kontinuerligt varierande och i {\"o}vrigt sn{\"a}ll funktion
verkligen skulle ge noll {\"o}verallt, eller hur?
Att funktionen har en singularitet i origo (placeringen av v{\aa}r
punktladdning) kan vi ignorera s{\aa} l{\"a}nge som vi utvecklar v{\aa}r
``derivata'' utanf{\"o}r singulariteten, s{\aa} faktum kvarst{\aa}r:
Hur kan en ``derivata'' av  $1/|{\bf x}|^2$ {\"o}verallt bli noll?

F{\"o}r den som fortfarande tvivlar p{\aa} att divergensen (v{\aa}r
``derivata'') f{\"o}r det elektriska f{\"a}ltet fr{\aa}n en punktladdning
{\"a}r noll, l{\aa}t oss derivera den spatiala variationen f{\"o}r det
ut{\aa}triktade f{\"a}ltet, som har formen
$$
  \eqalign{
   {\bf a}={\bf x}/|{\bf x}|^3\cr
    \quad\Rightarrow\quad
    \nabla\cdot{\bf a}
    &\equiv\Big(
        {{\partial}\over{\partial x}},
        {{\partial}\over{\partial y}},
        {{\partial}\over{\partial z}}
     \Big)\cdot{{(x,y,z)}\over{(x^2+y^2+z^2)^{3/2}}}\cr
    &={{1}\over{(x^2+y^2+z^2)^{3/2}}}
    \Big(
        {{\partial}\over{\partial x}},
        {{\partial}\over{\partial y}},
        {{\partial}\over{\partial z}}
     \Big)\cdot(x,y,z)
    +\Big(
        x{{\partial}\over{\partial x}},
        y{{\partial}\over{\partial y}},
        z{{\partial}\over{\partial z}}
     \Big)\cdot{{(1,1,1)}\over{(x^2+y^2+z^2)^{3/2}}}\cr
    &={{(1+1+1)}\over{(x^2+y^2+z^2)^{3/2}}}
        -{{3}\over{2}}(x,y,z)\cdot{{(2x,2y,2z)}\over{(x^2+y^2+z^3)^{5/2}}}\cr
    &={{3}\over{(x^2+y^2+z^3)^{3/2}}}
        -{{3}\over{2}}{{(2x^2+2y^2+2z^2)}\over{(x^2+y^2+z^3)^{5/2}}}\cr
    &={{3}\over{(x^2+y^2+z^3)^{3/2}}}
        -{{3}\over{(x^2+y^2+z^3)^{3/2}}}\cr
    &=0,\cr
  }
$$
det vill s{\"a}ga att divergensen {\it {\"o}verallt {\"a}r noll}, med undantag
f{\"o}r origo ${\bf x}={\bf 0}$ d{\"a}r divergensen {\"a}r odefinierad d{\aa}
n{\"a}mnaren d{\"a}r blir noll. Detta fall motsvarar det elektrostatiska
f{\"a}ltet fr{\aa}n en punktladdning placerad i origo.

\subsection{{\hl{En paradox f{\"o}r det elektrostatiska f{\"a}ltet i en homogen
  laddningst{\"a}thet}}}
Den f{\"o}reg{\aa}ende diskussionen kan l{\"a}tt ge intryck av att divergensen
av den elektriska f{\"a}ltet {\"a}r noll praktiskt taget {\"o}verallt, och att
vi i princip endast har att g{\"o}ra med delta-distributioner som k{\"a}llor.
Vi kommer d{\"a}rf{\"o}r att nu illustrera en situation d{\"a}r hela
dom{\"a}nen {\"a}r homogent fylld med elektrisk laddning, med andra ord att
vi har {\it k{\"a}llor {\"o}verallt}.
\epsfig{figs/field-unicharge.1}
\noindent
F{\"o}r fallet d{\aa} hela dom{\"a}nen {\"a}r homogent fylld med elektrisk
laddning har vi laddningst{\"a}theten $\rho({\bf x})=\rho_0=\hbox{konstant}$,
och divergensen f{\"o}r det statiska elektriska f{\"a}ltet {\"a}r d{\"a}rmed
{\"o}verallt noll-skild, som
$$
  \hl{\nabla\cdot{\bf E}({\bf x})=\rho_0/\varepsilon_0}.
$$
Samtidigt s{\aa} inser vi att vi i detta fall {\hl{av symmetrisk{\"a}l
m{\aa}ste ha att {\it det statiska elektriska f{\"a}ltet {\"o}verallt i
dom{\"a}nen med homogen laddning m{\aa}ste vara noll}}},
$$
  \hl{{\bf E}({\bf x})={\bf 0}}.
$$
{\it Man inser dock direkt att detta ser mycket m{\"a}rkligt ut, d{\aa} en
konstant funktion, och allra helst en funktion som {\"a}r identiskt noll,
inte kan ha en nollskild derivata.} Kort och gott, {\hl{hur g{\aa}r denna
paradox ihop?}} Det tycks ju onekligen vara s{\aa} att detta konstruerade (och
synnerligen enkla) exempel mots{\"a}ger Gauss lag f{\"o}r elektriska f{\"a}lt.

Svaret ligger i att vi m{\aa}ste vara f{\"o}rsiktiga i att anv{\"a}nda
symmetriargumentet, d{\"a}r vi f{\"o}r en {\"a}ndlig och symmetrisk
laddningsf{\"o}rdelning f{\"o}rvisso kan ha ett f{\"a}lt som {\"a}r noll i
centrum, men d{\"a}r f{\"a}ltet inte kan vara noll {\"o}verallt.
F{\"o}r en o{\"a}ndlig laddningsf{\"o}rdelning, som problemformuleringen ovan
implicit antar, uppst{\aa}r problem med randvillkor och fundamentet f{\"o}r
det elektriska f{\"a}ltets definition. Med andra ord kan vi inte anv{\"a}nda
symmetriargumentet f{\"o}r en {\"a}ndlig dom{\"a}n.

Sant {\"a}r dock att f{\"o}r en homogen laddningst{\"a}thet $\rho({\bf x})=
\rho_0$ s{\aa} {\"a}r alltid divergensen nollskild och beskriven av Gauss lag
$\nabla\cdot{\bf E}({\bf x})=\rho_0/\varepsilon_0$.
Om vi i en {\it {\"a}ndlig} och symmetrisk dom{\"a}n, exempelvis en homogent
laddad sf{\"a}r, betraktar det elektriska f{\"a}ltet, s{\aa} b{\"o}r
naturligtvis det elektriska f{\"a}ltet i centrum av sf{\"a}ren vara exakt noll,
men detta inneb{\"a}r ingen paradox, d{\aa} vi kan se detta helt analogt med,
s{\"a}g, en linj{\"a}r funktions nollgenomg{\aa}ng $f(x)=0$, med en konstant
derivata $df/dx=\hbox{konstant}$.

{\hl{\it Vi kan se just detta problem som ett typexempel p{\aa} en situation
d{\"a}r idealiserad fysikalisk situation inte alltid har en l{\"o}sning
som uppfyller alla antaganden!}}

\section{{\hl{Rotationen $\nabla\times{\bf b}$}}}
\subsection{Vad beskriver rotationen?}
Vi har sett att divergensen beskriver nettofl{\"o}det ut fr{\aa}n eller in mot
en punkt. {\it Rotationen} handlar ist{\"a}llet om hur ett vektorf{\"a}lt
lokalt {\it cirkulerar}.

F{\"o}rest{\"a}ll dig att du placerar ett litet, fritt roterande hjul eller en
paddel i en str{\"o}mmande v{\"a}tska. Om str{\"o}mningen f{\"o}rs{\"o}ker
vrida hjulet beskriver vi helt naturligt fl{\"o}det som {\it roterande i den
punkten}.
Denna analogi ger en intuitiv bild av rotationen, men vi m{\aa}ste h{\"a}r
skilja mellan ett vektorf{\"a}lts rotation och en enskild partikels r{\"o}relse.
Viktigt att notera {\"a}r att {\it ett vektorf{\"a}lt kan ha en noll-skild
rotation {\"a}ven om en partikel som f{\"o}ljer f{\"a}ltet inte r{\"o}r sig
l{\"a}ngs en cirkel.}

F{\"o}r ett vektorf{\"a}lt ${\bf b}$ definieras rotationen som
$$
  \hl{{\rm rot}({\bf b})\equiv\nabla\times{\bf b}},
$$
vilket vi i kartesiska koordinater kan uttrycka som
$$
  \nabla\times{\bf b}
    =\left|\matrix{
       {\bf e}_x & {\bf e}_y & {\bf e}_z\cr
       {{\partial}\over{\partial x}}
           & {{\partial}\over{\partial y}}
           & {{\partial}\over{\partial z}}\cr
       b_x & b_y & b_z\cr
     }\right|,
$$
eller utvecklat i komponenter,
$$
  \nabla\times{\bf b}
    ={\bf e}_x\Big(
      {{\partial b_z}\over{\partial y}}-{{\partial b_y}\over{\partial z}}
     \Big)
    +{\bf e}_y\Big(
      {{\partial b_x}\over{\partial z}}-{{\partial b_z}\over{\partial x}}
     \Big)
    +{\bf e}_z\Big(
      {{\partial b_y}\over{\partial x}}-{{\partial b_x}\over{\partial y}}
     \Big)
$$
Rotationen {\"a}r allts{\aa} i sig ett vektorf{\"a}lt. F{\"a}ltets riktning
anger rotationsaxeln enligt den vanliga h{\"o}gerhandsregeln, och dess belopp
anger den lokala rotationens styrka.

\subsection{Cirkulation och Stokes sats}
Den lokala rotationen {\"a}r kopplad till {\it cirkulationen av ett f{\"a}lt}
l{\"a}ngs en sluten kurva $\Gamma$, som definieras som
$$
  \oint_{\Gamma}{\bf b}\cdot d{\bf l}.
$$
Stokes teorem s{\"a}ger att cirkulationen l{\"a}ngs en sluten kurva lika med
fl{\"o}det av rotationen genom den yta som kurvan omsluter, vilket uttryckt
i vektoranalysens spr{\aa}k lyder
$$
  \oint_{\Gamma}{\bf b}\cdot d{\bf l}=\iint_S(\nabla\times{\bf b})\cdot d{\bf S},
$$
d{\"a}r $\Gamma$ {\"a}r randen (en sluten kurva) till ytan $S$, med orientering
av ytelementets riktning $d{\bf S}$ f{\"o}ljande h{\"o}gerhandsregeln.
Om rotationen $\nabla\times{\bf b}$ {\"a}r noll {\"o}verallt p{\aa} en
s{\aa}dan yta, blir ocks{\aa} cirkulationen noll l{\"a}ngs dess rand.

Det {\"a}r h{\"a}r viktigt att notera att rotationen inte {\"a}r samma sak som
att f{\"a}ltlinjerna visuellt ``ser b{\"o}jda ut''.
Ett f{\"a}lt kan ha kr{\"o}kta f{\"a}ltlinjer utan att ha n{\aa}gon rotation,
och omv{\"a}nt s{\aa} kan parallella f{\"a}ltlinjer fortfarande definiera ett
f{\"a}lt som har en nollskild rotation. Det avg{\"o}rande {\"a}r hur
f{\"a}ltets riktning och storlek varierar lokalt.

\subsection{Rotation hos det elektrostatiska f{\"a}ltet}
I elektro{\it statik} {\"a}r det elektriska f{\"a}ltet tidsoberoende och skapas
av station{\"a}ra laddningar. F{\"o}r detta station{\"a}ra fall g{\"a}ller
Faradays induktionslag i dess elektrostatiska form
$$
  \nabla\times{\bf E}={\bf 0}.
$$
Det elektro{\it statiska}\numberedfootnote{Notera att detta allts{\aa}
  g{\"a}ller ett elektro{\it statiskt} (tidsoberoende) f{\"a}lt; f{\"o}r ett
  elektro{\it dynamiskt} (tidsberoende) f{\"a}lt g{\"a}ller Faradays
  generella induktionslag
  $$
    \nabla\times{\bf E}=-{{\partial{\bf B}}\over{\partial t}},
  $$
  som vi kommer att studera i detalj och h{\"a}rleda i F{\"o}rel{\"a}sning~5.}
f{\"a}ltet {\"a}r allts{\aa} rotationsfritt. Detta betyder att linjeintegralen
av f{\"a}ltet mellan tv{\aa} punkter inte {\"a}r beroende av vilken v{\"a}g vi
v{\"a}ljer att utf{\"o}ra integralen l{\"a}ngs, f{\"o}rutsatt att vi r{\"o}r
oss inom ett l{\"a}mpligt omr{\aa}de d{\"a}r potentialen {\"a}r
v{\"a}ldefinierad.

Det {\"a}r just d{\"a}rf{\"o}r vi kan skriva det elektrostatiska f{\"a}ltet
som den negativa gradienten av den skal{\"a}ra elektrostatiska potentialen som
$$
  {\bf E}=-\nabla\phi.
$$
Vi erinrar oss att det {\it negativa tecknet f{\"o}r gradienten bara {\"a}r en
konvention} som vi genomg{\aa}ende anv{\"a}nder oss av, som s{\"a}ger att det
elektriska f{\"a}ltet definieras s{\aa} att f{\"a}ltlinjerna riktas fr{\aa}n
h{\"o}gre till l{\"a}gre potential.
H{\"a}r noterar vi att rotationen $\nabla\times{\bf E}$ {\"a}r noll {\"a}ven
i omr{\aa}den d{\"a}r det finns elektriska laddningar, s{\aa} l{\"a}nge vi
betraktar det elektrostatiska f{\"a}ltet p{\aa} ett s{\"a}tt som g{\"o}r det
definierat\numberedfootnote{Det vill s{\"a}ga att vi exkluderar eventuella
  singulariteter och gr{\"a}ns{\"o}verg{\aa}ngar mellan dom{\"a}ner med,
  s{\"a}g, olika elektrisk permittivitet.}
Vi kan konstatera att {\it ger laddningar i vila upphov till divergens, men
aldrig rotation}, n{\aa}got som direkt f{\"o}ljer av vektoridentiteten
$$
  \nabla\times(\nabla\phi)\equiv{\bf 0}.
$$
Det {\"a}r endast d{\aa} elektriska laddningar s{\"a}tts i r{\"o}relse och
d{\"a}rmed genererar en str{\"o}m som de kan ge upphov till en nollskild
rotation f{\"o}r det elektriska f{\"a}ltet, och d{\aa} via Faradays lag
enligt ovan.

\subsection{Rotation hos det magnetostatiska f{\"a}ltet}
F{\"o}r magnetf{\"a}ltet {\"a}r situationen annorlunda, och vi har f{\"o}r
dessa f{\"a}lt Amp{\`e}res lag\numberedfootnote{Eller om man s{\aa} vill,
{\it Amp{\`e}re--Maxwells lag}} f{\"o}r vakuum som
$$
  \nabla\times{\bf B}
    =\mu_0{\bf J}+\mu_0\varepsilon_0{{\partial{\bf E}}\over{\partial t}}
$$
d{\"a}r ${\bf J}$ {\"a}r str{\"o}mt{\"a}theten (${\rm A}/{\rm m}^2$).
I magnetostatik {\"a}r f{\"a}lten tidsoberoende, s{\aa} att
$\partial{\bf E}/\partial t=0$, och vi f{\aa}r f{\"o}r det
magneto{\it statiska} fallet att
$$
  \nabla\times{\bf B}=\mu_0{\bf J}.
$$
Detta inneb{\"a}r att det statiska magnetf{\"a}ltet i allm{\"a}nhet inte {\"a}r
rotationsfritt, eftersom n{\"a}rvaro av elektrisk str{\"o}m ger att rotationen
av ${\bf B}$ i allm{\"a}nhet {\"a}r skild fr{\aa}n noll.

I ett {\it str{\"o}mfritt omr{\aa}de} g{\"a}ller det d{\"a}remot alltid att
$$
  \nabla\times{\bf B}=0.
$$
Observera att detta endast s{\"a}ger att f{\"a}ltet {\"a}r rotationsfritt i det
aktuella omr{\aa}det, och inneb{\"a}r inte n{\"o}dv{\"a}ndigtvis att f{\"a}ltet
{\"a}r noll d{\"a}r. Ett magnetf{\"a}lt kan finnas i ett omr{\aa}de d{\"a}r
ingen str{\"o}m flyter, samtidigt som dess rotation {\"a}r noll i det
omr{\aa}det.

Ett klassiskt exempel {\"a}r omr{\aa}det utanf{\"o}r en l{\aa}ng, rak ledare
d{\"a}r magnetf{\"a}ltet {\"a}r riktat cirkul{\"a}rt runt ledaren.
F{\"a}ltet har i detta fall {\"a}nd{\aa} noll rotation i varje punkt
utanf{\"o}r ledaren, men den slutna linjeintegralen runt ledaren kan
{\"a}nd{\aa} vara skild fr{\aa}n noll eftersom det omr{\aa}de som den
omslutande kurvan sp{\"a}nner upp inneh{\aa}ller str{\"o}mmen och inte
{\"a}r ett helt str{\"o}mfritt, enkelt sammanh{\"a}ngande omr{\aa}de.

Detta exempel visar varf{\"o}r man m{\aa}ste vara f{\"o}rsiktig n{\"a}r man
anv{\"a}nder lokala differentialegenskaper f{\"o}r att tolka och dra slutsatser
om integraler l{\"a}ngs slutna kurvor.

\subsection{{\hl{Biot--Savarts lag och rotationen hos magnetf{\"a}ltet}}}
Enligt Biot–Savarts lag ges magnetf{\"a}ltet fr{\aa}n en o{\"a}ndligt l{\aa}ng,
rak str{\"o}mf{\"o}rande ledare, b{\"a}rande str{\"o}mmen $I$ l{\"a}ngs
$z$-axeln genom origo, som
$$
  \hl{{\bf B}({\bf x})={{\mu_0I}\over{2\pi r}}{\bf e}_{\varphi}},
$$
d{\"a}r {\hl{$r=(x^2+y^2)^{1/2}$}} {\"a}r det vinkelr{\"a}ta avst{\aa}ndet till
$z$-axeln och ${\bf e}_{\varphi}$ {\"a}r den pol{\"a}ra riktningen i cylindriska
koordinater.
\epsfig{figs/field-biot.1}
\noindent
I kartesiska koordinater har vi detta uttryckt som
$$
  \hl{{\bf e}_{\varphi}=
    -{\bf e}_x{{y}\over{(x^2+y^2)^{1/2}}}
    +{\bf e}_x{{x}\over{(x^2+y^2)^{1/2}}}},
$$
och fr{\aa}n dessa identiteter erh{\aa}ller vi d{\"a}rmed att
$$
  \hl{{\bf B}({\bf x})={{\mu_0I}\over{2\pi(x^2+y^2)}}
    \big(-y{\bf e}_x+x{\bf e}_y\big)}.
$$
Rotationen av magnetf{\"a}ltet f{\aa}r vi d{\aa} som
$$
  \eqalign{
    \hl{\nabla\times{\bf B}({\bf x})}&\hl{={{\mu_0I}\over{2\pi}}
    \bigg(
      {\bf e}_x{{\partial}\over{\partial x}}
        +{\bf e}_y{{\partial}\over{\partial y}}
        +\underbrace{{\bf e}_z{{\partial}\over{\partial z}}}_{\to0}
    \bigg)
    \times
    \bigg(
      -{{y}\over{(x^2+y^2)}}{\bf e}_x+{{x}\over{(x^2+y^2)}}{\bf e}_y
    \bigg)}\cr
    &\hl{={{\mu_0I}\over{2\pi}}
    {\bf e}_z
    \bigg(
      {{\partial}\over{\partial x}}\bigg({{x}\over{x^2+y^2}}\bigg)
        +{{\partial}\over{\partial y}}\bigg({{y}\over{x^2+y^2}}\bigg)
    \bigg)}\cr
    &\hl{={{\mu_0I}\over{2\pi}}
    {\bf e}_z
    \bigg(
      \bigg({{1}\over{x^2+y^2}}-{{x\cdot2x}\over{(x^2+y^2)^2}}\bigg)
        +\bigg({{1}\over{x^2+y^2}}-{{y\cdot2y}\over{(x^2+y^2)^2}}\bigg)
    \bigg)}\cr
    &\hl{={{\mu_0I}\over{2\pi}}
    {\bf e}_z\bigg({{2(x^2+y^2)-2(x^2+y^2)}\over{(x^2+y^2)^2}}\bigg)}\cr
    &\hl{={\bf 0}}.\cr
  }
$$
Med andra ord har vi en liknande situation som tidigare f{\"o}r divergensen,
n{\"a}mligen att rotationen {\"o}verallt {\"a}r noll, f{\"o}rutom i origo,
trots att det kontinuerliga uttrycket f{\"o}r magnetf{\"a}ltet onekligen ser
ut att ha en nollskild ``derivata''.

Det centrala {\"a}r h{\"a}r att divergensen och rotationen m{\"a}ter olika
saker. Divergensen beskriver nettofl{\"o}de ut fr{\aa}n eller in mot en punkt,
medan rotationen beskriver lokal cirkulation. Ett f{\"a}lt kan mycket v{\"a}l
ha noll av den ena men inte av den andra.

\subsection{{\hl{Kan parallella f{\"a}ltlinjer ge upphov till en rotation?}}}
Vid en f{\"o}rsta anblick verkar ju denna fr{\aa}ga ganska konstig.
Om rotationen i n{\aa}gon bem{\"a}rkelse {\"a}r ett m{\aa}tt p{\aa} hur mycket
f{\"a}ltet ``snurrar runt'' lokalt, s{\aa} m{\aa}ste v{\"a}l {\"a}nd{\aa}
f{\"a}ltlinjerna p{\aa} n{\aa}got vis kr{\"o}kas, eller har vi p{\aa} n{\aa}got
s{\"a}tt missf{\"o}rst{\aa}tt fr{\aa}gan?
\epsfig{figs/field-rot-parallel.1}
\noindent
F{\"o}r att illustrera po{\"a}ngen, l{\aa}t oss konstruera ett artificiellt
inhomogent vektorf{\"a}lt ${\bf a}({\bf x})$ riktat l{\"a}ngs $x$-axeln och
varierande i $y$-led enligt
$$
  \hl{{\bf a}({\bf x})=(a_0-a_1y){\bf e}_x}.
$$
F{\"o}r detta vektorf{\"a}lt, som helt uppenbarligen har alla f{\"a}ltlinjer
parallella med varandra, blir rotationen av f{\"a}ltet
$$
  \eqalign{
    \hl{\nabla\times{\bf a}({\bf x})}&=\hl{
    \bigg(
      \underbrace{{\bf e}_x{{\partial}\over{\partial x}}}_{\to0}
        +{\bf e}_y{{\partial}\over{\partial y}}
        +\underbrace{{\bf e}_z{{\partial}\over{\partial z}}}_{\to0}
    \bigg)\times\Big((a_0-a_1y){\bf e}_x\Big)}\cr
    &=\hl{\underbrace{{\bf e}_y\times{\bf e}_x}_{=-{\bf e}_z}
        {{\partial}\over{\partial y}}(a_0-a_1y)}\cr
    &=\hl{a_1{\bf e}_z},\cr
  }
$$
det vill s{\"a}ga att rotationen av f{\"a}ltet {\"o}verallt {\"a}r nollskild,
trots att samtliga f{\"a}ltlinjer {\"a}r parallella!

F{\"o}r att visualisera vad som f{\"o}rsigg{\aa}r h{\"a}r kan vi ta till
v{\aa}rt ``mentala paddel'' som vi t{\"a}nker oss att vi stoppar ner i en
inhomogent men parallellt flytande str{\"o}m av v{\"a}tska. Kommer v{\aa}r
``virtuella paddel'' att b{\"o}rja rotera i det inhomogena fl{\"o}det som
beskrivs i uppgiften?
Ja, det inser vi naturligtvis att den kommer att g{\"o}ra, och vi inser
f{\"o}rhoppningsvis att {\hl{\it en nollskild rotation inte
n{\"o}dv{\"a}ndigtvis beh{\"o}ver h{\"a}nga samman med att f{\"a}ltet
kr{\"o}ker sig}}.
\vfill\eject

\subsection{{\hl{Exempel p{\aa} f{\"a}lt med en homogen rotation}}}
Vi avslutar avsnittet om rotationen med att konstatera att det naturligtvis
inte bara {\"a}r fr{\aa}ga om att rotationen av ett magnetf{\"a}lt, enligt
Biot--Savarts lag, i princip alltid {\"a}r noll {\"o}verallt och att k{\"a}llor
till magnetf{\"a}lt i princip alltid skulle vara givna av delta-distributioner.
\epsfig{figs/field-unicurrent.1}
\noindent
Om vi exempelvis har en dom{\"a}n som genoml{\"o}ps av en homogen
str{\"o}mt{\"a}thet ${\bf J}({\bf x})=J_0{\bf e}_z$, s{\aa} s{\"a}ger
Amp{\`e}res statiska lag f{\"o}r magnetf{\"a}lt att rotationen av
magnetf{\"a}ltet {\"o}verallt {\"a}r konstant och riktat l{\"a}ngs
$z$-axeln, som
$$
  \hl{\nabla\times{\bf B}=\mu_0{\bf J}=\mu_0J_0{\bf e}_z}.
$$
Liksom tidigare, i analysen av paradoxen f{\"o}r divergensen
$\nabla\cdot{\bf E}$ av ett elektrostatiskt f{\"a}lt i fallet med en homogen
laddningst{\"a}thet, s{\aa} tycks vi h{\"a}r av symmetrisk{\"a}l ha att
magnetf{\"a}ltet {\"o}verallt {\"a}r noll, d{\aa} ``intilliggande virvlar
av f{\"a}ltlinjer'' ju {\"o}verallt kancellerar varandra.
{\AA}terigen, s{\aa} ligger l{\"o}sning f{\"o}r denna paradox i att vi inte kan
anta en o{\"a}ndlig dom{\"a}n utan att bryta mot grundl{\"a}ggande randvillkor,
denna g{\aa}ng f{\"o}r magnetf{\"a}ltets definition.

{\hl{\it {\AA}terigen kan vi se detta problem som ett typexempel p{\aa} en
situation d{\"a}r idealiserad fysikalisk situation inte alltid har en
l{\"o}sning som uppfyller alla antaganden!}}

\subsection{Kan en ytstr{\"o}m ge rotation skild fr{\aa}n noll?}
En ytstr{\"o}m {\"a}r ett annat intressant exempel som anknyter till v{\aa}r
f{\"o}rsta diskussion kring Diracs delta-distribution, eftersom en ytstr{\"o}m
{\"a}r koncentrerad till en yta snarare {\"a}n som i f{\"o}reg{\aa}ende exempel
utspridd {\"o}ver en volym.

Antag att vi har en o{\"a}ndlig str{\"o}mskiv i planet $z=0$, med
ytstr{\"o}mt{\"a}theten
$$
  {\bf K}=K_0{\bf e}_x,
$$
d{\"a}r ytstr{\"o}mt{\"a}theten ${\bf K}$ har enheten ${\rm A}/{\rm m}$,
f{\"o}r vilken den motsvarande volymstr{\"o}mt{\"a}theten kan uttryckas med
en deltafunktion som
$$
  {\bf J}=K_0\delta(z){\bf e}_x.
$$
Amp{\`e}res statiska lag ger d{\aa} direkt att
$$
  \nabla\times{\bf B}({\bf x})=\mu_0{\bf J}=\mu_0K_0\delta(z){\bf e}_x.
$$
Rotationen {\"a}r noll p{\aa} vardera sidan om skivan, men har ett
infinitesimalt utstr{\"a}ckt bidrag fr{\aa}n delta-distributionen i $z$-planet,
vilket beskriver hur magnetf{\"a}ltet ${\bf B}({\bf x})$ {\"a}ndrar sig
spr{\aa}ngvis {\"o}ver ytstr{\"o}mmen. Detta kan vi se som ett exempel p{\aa}
skillnaden mellan en rotation som {\"a}r skild fr{\aa}n noll i ett helt
volymomr{\aa}de och en rotation som {\"a}r koncentrerad till en idealiserad
yta.
\vfill\eject

\section{Konservativa f{\"a}lt}
\subsection{Vad {\"a}r ett konservativt f{\"a}lt?}
Vi har nu analyserat tv{\aa} s{\"a}tt att beskriva ett vektorf{\"a}lts lokala
egenskaper, geonom divergens och rotation, och vi kommer nu att introducera
ett begrepp som knyter an till b{\aa}da dessa operatorer, n{\"a}mligen
{\it konservativa f{\"a}lt}.
\quote{%
  {\hl{{\bf Definition:} Ett vektorf{\"a}lt ${\bf F}$ {\"a}r {\it konservativt}
  i ett omr{\aa}de om dess {\it linjeintegral mellan tv{\aa} punkter endast
  beror p{\aa} startpunkten och slutpunkten}, och inte p{\aa} vilken v{\"a}g
  vi tar d{\"a}remellan.}}
}
\noindent
L{\aa}t oss beteckna startpunkten med ${\bf x}_A$ och slutpunkten med
${\bf x}_B$. F{\"o}r ett konservativt f{\"a}lt g{\"a}ller d{\aa} att
$$
  \int^{{\bf x}_B}_{{\bf x}_A}{\bf F}\cdot d{\bf l}=\phi({\bf x}_A)-\phi({\bf x}_B)
$$
f{\"o}r n{\aa}gon skal{\"a}r potential $\phi$, om vi v{\"a}ljer
konventionen\numberedfootnote{Det {\"a}r h{\"a}r viktigt att h{\aa}lla
  i minnet vilken konvention f{\"o}r potential som vi anv{\"a}nder.
  Vi anv{\"a}nder ju h{\"a}r den elektrostatiska definitionen med ett
  minustecken framf{\"o}r gradienten, och detta {\"a}r anledningen till
  att ordningen p{\aa} termerna i h{\"o}gerledet kan se lite udda ut, om
  vi j{\"a}mf{\"o}r med den rent visuellt liknande {\it integralkalkylens
  huvudsats} ({\"a}ven kallad {\it analysens fundamentalsats}),
  $$
    \int^b_a f(x)\,dx=F(b)-F(a).
  $$}
$$
  {\bf F}=-\nabla\phi.
$$
Det finns allts{\aa} en skal{\"a}r potential vars negativa gradient ger det
konservativa f{\"a}ltet.

En ekvivalent egenskap {\"a}r att linjeintegralen l{\"a}ngs varje sluten kurva
{\"a}r noll,
$$
  \oint_{\Gamma}{\bf F}\cdot d{\bf l}=0,
$$
vilket {\"a}r ett mycket anv{\"a}ndbart kriterium f{\"o}r konservativa f{\"a}lt.
Om vi g{\aa}r runt en sluten slinga $\Gamma$ och kommer tillbaka till
utg{\aa}ngspunkten kan ett konservativt f{\"a}lt inte ge n{\aa}gon
nettointegral l{\"a}ngs slingan. J{\"a}mf{\"o}r g{\"a}rna med v{\aa}r mentala
bild av en skal{\"a}r potential som varande en ``h{\"o}jd i ett landskap'',
enligt tidigare.

S{\aa} varf{\"o}r {\"a}r d{\aa} konservativa f{\"a}lt s{\aa} fundamentalt
intressanta inom elektromagnetism? Svaret p{\aa} denna fr{\aa}ga {\"a}r i
grund och botten att konservativa f{\"a}lt per automatik inneb{\"a}r att vi
kan anv{\"a}nda potentialen f{\"o}r att f{\"o}renkla ber{\"a}kningar.
Ist{\"a}llet f{\"o}r att integrera f{\"a}ltet l{\"a}ngs en komplicerad v{\"a}g
kan vi best{\"a}mma potentialen i tv{\aa} punkter och ta skillnaden, helt i
analogi med klassiska mekaniska problem d{\"a}r vi ofta kan f{\"o}renkla
problem om vi kan identifiera en potential.

\subsection{Exempel fr{\aa}n elektrostatik}
Det elektrostatiska f{\"a}ltet {\"a}r ett klassiskt exempel p{\aa} ett
konservativt f{\"a}lt, f{\"o}r att inte s{\"a}ga {\it det} klassiska
exemplet. Eftersom vi f{\"o}r ett elektro\-statiskt f{\"a}lt har att
$$
  {\bf E}=-\nabla\phi,
$$
s{\aa} f{\aa}r vi med hj{\"a}lp av linjeintegralen f{\"o}r gradienten att
$$
  \int^{{\bf x}_B}_{{\bf x}_A}{\bf E}\cdot d{\bf l}
    =-\int^{{\bf x}_B}_{{\bf x}_A}\nabla\phi\cdot d{\bf l}.
$$
Integralen av en gradient l{\"a}ngs en kurva {\"a}r potentialens
f{\"o}r{\"a}ndring mellan {\"a}ndpunkterna, och d{\"a}rf{\"o}r g{\"a}ller
det att
$$
  \int^{{\bf x}_B}_{{\bf x}_A}{\bf E}\cdot d{\bf l}=\phi({\bf x}_A)-\phi({\bf x}_B).
$$
Den elektriska potentialskillnaden (volt) mellan punkterna ${\bf x}_A$ och
${\bf x}_B$ {\"a}r s{\aa}ledes oberoende av vilken v{\"a}g vi v{\"a}ljer i
ett elektrostatiskt f{\"a}lt. Denna observation utg{\"o}r i sj{\"a}lva verket
grunden f{\"o}r begreppet {\it sp{\"a}nning}, som kort och gott definieras som
{\it potentialskillnaden mellan tv{\aa} punkter}.

F{\"o}r en positiv testladdning $q$ {\"a}r arbetet som det elektriska
f{\"a}ltet utf{\"o}r n{\"a}r laddningen f{\"o}rflyttas fr{\aa}n ${\bf x}_A$
till ${\bf x}_B$
$$
  \Delta W({\bf x}_A\to{\bf x}_B)=q\int^{{\bf x}_B}_{{\bf x}_A}
    {\bf E}\cdot d{\bf l} =q\big(\phi({\bf x}_A)-\phi({\bf x}_A)\big).
$$
Arbetet beror allts{\aa} endast p{\aa} {\"a}ndpunkterna, och om laddningen
f{\"o}rs tillbaka till utg{\aa}ngspunkten utf{\"o}r f{\"a}ltet inget
nettoarbete. Detta {\"a}r den fysikaliska inneb{\"o}rden av att f{\"a}ltet
{\"a}r konservativt, i bem{\"a}rkelsen att den totala energin bevaras utan
att {\"o}verg{\aa} till v{\"a}rme.

\subsection{{\"A}r magnetf{\"a}ltet konservativt?}
H{\"a}r beh{\"o}ver vi g{\"o}ra en distinktion mellan magnetf{\"a}ltet%
\numberedfootnote{{\AA}terigen, om vi skall vara petiga, den
  {\it magnetiska fl{\"o}dest{\"a}theten} ${\bf B}$.}
${\bf B}$, den magnetiska kraften och den magnetiska vektorpotentialen
${\bf A}$. I magnetostatik g{\"a}ller Amp{\`e}res lag
$$
  \nabla\times{\bf B}=\mu_0{\bf J},
$$
och om str{\"o}mt{\"a}theten {\"a}r noll har vi d{\"a}rf{\"o}r
$$
  \nabla\times{\bf B}=0.
$$
Man kan tyckas att detta borde inneb{\"a}ra att magnetf{\"a}ltet {\"a}r
konservativt. Lokalt, i ett enkelt sammanh{\"a}ngande och str{\"o}mfritt
omr{\aa}de, kan vi faktiskt representera det statiska magnetiska f{\"a}ltet
som en gradient
$$
  {\bf B}=-\nabla\psi,
$$
d{\"a}r $\psi$ {\"a}r en skal{\"a}r potential f{\"o}r det statiska
magnetf{\"a}ltet. Detta samband g{\"a}ller dock {\it inte} generellt,
och vi kommer ej heller i kursen att anv{\"a}nda begreppet magnetisk potential.

Betrakta {\aa}terigen magnetf{\"a}ltet kring en l{\aa}ng, rak
str{\"o}mf{\"o}rande ledare, d{\"a}r Biot--Savarts lag,
$$
  {\bf B}({\bf x})={{\mu_0I}\over{2\pi r}}{\bf e}_{\varphi},
$$
s{\"a}ger oss att magnetf{\"a}ltet g{\aa}r i cirklar runt ledaren. En sluten
kurva $\Gamma$ som omsluter ledaren, d{\"a}r vi l{\"a}mpligen kan v{\"a}lja
den slutna kurvan som en cirkel centrerad vid ledaren och med radie $r$, ger
oss direkt att den slutna linjeintegralen runt denna blir
$$
  \oint_{\Gamma}{\bf B}({\bf x})\cdot d{\bf l}
    ={{\mu_0I}\over{2\pi r}}2\pi r
    =\mu_0 I.
$$
Om str{\"o}mmen $I$ {\"a}r skild fr{\aa}n noll blir integralen allts{\aa}
{\it inte} noll, och magnetf{\"a}ltet {\"a}r d{\"a}rf{\"o}r {\it inte
konservativt} i ett omr{\aa}de som inneh{\aa}ller en s{\aa}dan slinga och
d{\"a}r den omslutna str{\"o}mmen {\"a}r skild fr{\aa}n noll.

Det viktiga i detta resonemang {\"a}r att ett f{\"a}lt kan vara rotationsfritt
i varje punkt i ett visst omr{\aa}de, men {\"a}nd{\aa} sakna en globalt entydig
skal{\"a}rpotential om omr{\aa}det inte {\"a}r enkelt sammanh{\"a}ngande.

\subsection{Konservativa f{\"a}lt i elektrodynamik}
I elektro{\it dynamik} varierar f{\"a}lten med tiden, och Faradays
induktionslag ger d{\aa} vid hand att
$$
  \nabla\times{\bf E}=-{{\partial{\bf B}}\over{\partial t}}.
$$
Om magnetf{\"a}ltet varierar med tiden kan det elektriska f{\"a}ltet f{\aa} en
rotation som inte {\"a}r noll, och vi har d{\aa} i allm{\"a}nhet {\it inte} att
det elektriska f{\"a}ltet kan skrivas som enbart beroende av den skal{\"a}ra
potentialen med ${\bf E}=-\nabla\phi$ f{\"o}r hela det elektriska f{\"a}ltet.
I detta fall anv{\"a}nder vi ist{\"a}llet det fullst{\"a}ndiga sambandet%
\numberedfootnote{Denna form och kopplingen till den elektromagnetiska
  v{\aa}gekvationen g{\aa}s igenom i F{\"o}rel{\"a}sning~9.}
$$
  {\bf E}=-\nabla\phi-{{\partial{\bf A}}\over{\partial t}},
$$
vilket delar upp det elektriska f{\"a}ltet i en del fr{\aa}n den gradienten
av den skal{\"a}ra potentialen $\phi$ och en del fr{\aa}n tidsderivatan av
vektorpotentialen ${\bf A}$. Detta inneb{\"a}r att det totala elektrodynamiska
f{\"a}ltet ${\bf E}({\bf x},t)$ inte generellt {\"a}r konservativt; dock, om
det magnetiska f{\"a}ltet ${\bf B}$ {\"a}r tidsoberoende i det aktuella
omr{\aa}det, s{\aa} att
$$
  {{\partial{\bf B}}\over{\partial t}}=0,
$$
{\"a}r det elektriska f{\"a}ltet rotationsfritt d{\"a}r. Under l{\"a}mpliga
topologiska f{\"o}ruts{\"a}ttningar kan det d{\aa} beskrivas med en skal{\"a}r
potential.
\bigskip
\noindent
{\bf Tabell 1}. N{\aa}gra typfall f{\"o}r konservativa respektive
  icke-konservativa f{\"a}lt.
\smallskip
{
\newdimen\hcellpad
\newdimen\vcellpad
\hcellpad=4pt
\vcellpad=4pt
\offinterlineskip
%\centerline{%
%\noindent{{\bf Tabell 1.} Egenskaper hos ........}
\halign{%
  \vrule
  \hbox to 120pt{\hskip\hcellpad
    \vbox{\vskip\vcellpad\hbox to 100pt{\hss#\hss}\vskip\vcellpad}%
    \hskip\hcellpad}
  \vrule&
  \hbox to 86pt{\hskip\hcellpad
    \vbox{\vskip\vcellpad\hbox to 78pt{\hss#\hss}\vskip\vcellpad}%
    \hskip\hcellpad}
  \vrule&
  \hbox to 220pt{\hskip\hcellpad
    \vbox{\vskip\vcellpad\hbox to 203pt{\hss#\hss}\vskip\vcellpad}%
    \hskip\hcellpad}
  \vrule\cr
  \noalign{\hrule}
  {\bf F{\"a}lt och situation}
    & {\bf Rotation}
    & {\bf Konservativt f{\"a}lt?} \cr
  \noalign{\hrule}
  Elektrostatiskt ${\bf E}$
    & $\nabla\times{\bf E}={\bf 0}$
    & Ja, i ett l{\"a}mpligt sammanh{\"a}ngande omr{\aa}de\cr
  \noalign{\hrule}
  \vbox{\hbox{Elektrodynamiskt ${\bf E}$,}\hbox{varierande ${\bf B}$}}
    & $\displaystyle\nabla\times{\bf E}=-{{\partial{\bf B}}\over{\partial t}}$
    & Inte generellt \cr
  \noalign{\hrule}
  \vbox{\hbox{Magnetostatiskt ${\bf B}$,}\hbox{str{\"o}mfritt omr{\aa}de}}
    & $\nabla\times{\bf B}={\bf 0}$
    & Ja lokalt, om omr{\aa}det har l{\"a}mplig topologi \cr
  \noalign{\hrule}
  \vbox{\hbox{Magnetostatiskt ${\bf B}$,}\hbox{omr{\aa}de med str{\"o}m}}
    & $\nabla\times{\bf B}=\mu_0{\bf J}$
    & Inte generellt \cr
  \noalign{\hrule}
}
%}
}
\bigskip

\subsection{Under vilka villkor {\"a}r ett f{\"a}lt konservativt?}
Det finns ett viktigt matematiskt samband mellan konservativa f{\"a}lt och
rotationsfria f{\"a}lt. Om ett f{\"a}lt kan skrivas som en gradient av
n{\aa}gon skal{\"a}r potential,
$$
  {\bf F}=\nabla\psi,
$$
s{\aa} {\"a}r f{\"a}ltet {\it alltid rotationsfritt} som en f{\"o}ljd av
vektoridentiteten\numberedfootnote{Detsamma g{\"a}ller sj{\"a}lvfallet {\"a}ven
  om vi anv{\"a}nder ``minus-konventionen'' (som inom elektrostatiken) och
  skriver ${\bf F}=-\nabla\phi$.}
$$
  \nabla\times(\nabla\psi)=0.
$$
Omv{\"a}nt g{\"a}ller att ett kontinuerligt och tillr{\"a}ckligt glatt f{\"a}lt
med noll rotation kan skrivas som en gradient i ett enkelt sammanh{\"a}ngande
omr{\aa}de. Ett enkelt sammanh{\"a}ngande omr{\aa}de {\"a}r, informellt, ett
omr{\aa}de utan h{\aa}l som hindrar oss fr{\aa}n att kontinuerligt dra ihop
varje sluten kurva till en punkt.

Med andra ord,
$$
  \nabla\times{\bf F}=0 \qquad\Rightarrow\qquad {\bf F}=-\nabla\phi
$$
under l{\"a}mpliga f{\"o}ruts{\"a}ttningar p{\aa} omr{\aa}det som f{\"a}ltet
${\bf F}$ {\"a}r definierat p{\aa}. Noll rotation r{\"a}cker inte alltid
f{\"o}r att garantera att det existerar en global potential $\phi$ om
omr{\aa}det har h{\aa}l eller andra topologiska hinder.

\section{{\hl{Tolkningen av $\nabla\times(\nabla\phi)$}}}
Vi kommer nu att g{\aa} igenom tre vektoridentiteter som {\"a}r av fundamentalt
intresse inom elektromagnetiska problem. Den f{\"o}rsta identiteten s{\"a}ger
oss att {\it rotationen av en gradient alltid {\"a}r identiskt noll,}
$$
  \nabla\times(\nabla\phi)={\bf 0}.
$$\sidx{Vektoridentitet}[$\nabla\times(\nabla\phi)={\bf 0}$]
F{\"o}r att illustrera vad denna vektoridentitet betyder, b{\"o}rjar vi med en
skal{\"a}r funktion $\phi$. D{\aa} vi tar gradienten av denna funktion, eller
skal{\"a}ra {\it f{\"a}lt}, f{\"o}r att anknyta till b{\"o}rjan av denna extra
f{\"o}rel{\"a}sning, s{\aa} omvandlar vi den till ett vektorf{\"a}lt
$\phi\rightarrow\nabla\phi$. Om vi d{\"a}refter tar rotationen av detta
vektorf{\"a}lt f{\aa}r vi alltid noll, f{\"o}rutsatt att $\phi$ {\"a}r
``tillr{\"a}ckligt glatt'' s{\aa} att de blandade andraderivatorna {\"a}r lika
n{\"a}r vi byter ordning f{\"o}r differentieringen.

F{\"o}r att se varf{\"o}r kan vi skriva ut rotationen i komponenter, d{\"a}r
exempelvis $z$-komponenten {\"a}r given som
$$
   (\nabla\times\nabla\phi)_z
      ={{\partial}\over{\partial x}}
          \bigg({{\partial\phi}\over{\partial y}}\bigg)
          -{{\partial}\over{\partial y}}
               \bigg({{\partial\phi}\over{\partial x}}\bigg)
      ={{\partial^2\phi}\over{\partial x\,\partial y}}
        -{{\partial^2\phi}\over{\partial y\,\partial x}}
      =0.
$$
Den sista likheten f{\"o}ljer av att blandade andraderivator {\"a}r lika
f{\"o}r en ``tillr{\"a}ckligt glatt'' funktion, och alla {\"o}vriga komponenter
blir noll p{\aa} samma s{\"a}tt. {\it Den geometriska tolkningen {\"a}r att
ett gradientf{\"a}lt aldrig kan ha n{\aa}gon lokal rotation}, vilket i sin tur
f{\"o}rklarar varf{\"o}r elektrostatiska f{\"a}lt, som kan skrivas som
$-\nabla\phi$, alltid {\"a}r rotationsfria.

Vektoridentiteten $\nabla\times(\nabla\phi)={\bf 0}$ f{\"o}rklarar ocks{\aa}
varf{\"o}r vi kan g{\"o}ra en s{\aa} kallad {\it gaugetransformation} av
vektorpotentialen som
$$
  {\bf A}'={\bf A}+\nabla\chi,
$$
detta eftersom magnetf{\"a}ltet ${\bf B}=\nabla\times{\bf A}$ definieras genom
rotationen av vektorpotentialen, med f{\"o}ljden att f{\"a}ltet inte
p{\aa}verkas av en transformation d{\"a}r vi l{\"a}gger till en gradient,
$$
  \nabla\times{\bf A}'
    =\nabla\times{\bf A}+\underbrace{\nabla\times(\nabla\chi)}_{=0}
    =\nabla\times{\bf A}.
$$
Gradienten av en deriverbar men i {\"o}vrigt godtycklig skal{\"a}r funktion
$\chi$ kan med andra ord f{\"o}r{\"a}ndra vektorpotentialen utan att
f{\"o}r{\"a}ndra det magnetf{\"a}lt som den representerar. Vi g{\aa}r igenom
gaugetransformationen och hur den {\"a}r anv{\"a}ndbar f{\"o}r att ta fram
v{\aa}gekvationer f{\"o}r den s{\aa} kallade {\it retarderade potentialen}
i F{\"o}rel{\"a}sning~11.

\section{{\hl{Tolkningen av $\nabla\cdot(\nabla\times{\bf a})$}}}
Den andra identiteten s{\"a}ger oss att {\it divergensen av en rotation alltid
{\"a}r identiskt noll,}
$$
  \nabla\cdot(\nabla\times{\bf a})=0.
$$\sidx{Vektoridentitet}[$\nabla\cdot(\nabla\times{\bf a})=0$]
Vi tar f{\"o}rst rotationen av vektorf{\"a}ltet ${\bf a}$ och f{\aa}r ett
nytt vektorf{\"a}lt ${\bf a}\rightarrow\nabla\times{\bf a}$, varefter vi tar
divergensen av detta f{\"a}lt, med resultatet att denna divergens alltid {\"a}r
identiskt noll om ${\bf a}$ {\"a}r ``tillr{\"a}ckligt glatt''.

F{\"o}r att se hur denna identitet uppkommer, kan vi utveckla divergensen av
rotationen av ett generellt vekt{\"o}rf{\"a}lt ${\bf a}$ som
$$
  \eqalign{
    \nabla\cdot(\nabla\times{\bf a})
      &={{\partial}\over{\partial x}}
        \bigg(
          {{\partial a_z}\over{\partial y}}-{{\partial a_y}\over{\partial z}}
        \bigg)
       +{{\partial}\over{\partial y}}
        \bigg(
          {{\partial a_x}\over{\partial z}}-{{\partial a_z}\over{\partial x}}
        \bigg)
       +{{\partial}\over{\partial z}}
        \bigg(
          {{\partial a_y}\over{\partial x}}-{{\partial a_x}\over{\partial y}}
        \bigg)\cr
      &=\bigg(
          {{\partial^2 a_z}\over{\partial x\partial y}}
             -{{\partial^2 a_y}\over{\partial x\partial z}}
        \bigg)
       +\bigg(
          {{\partial^2 a_x}\over{\partial y\partial z}}
             -{{\partial^2 a_z}\over{\partial y\partial x}}
        \bigg)
       +\bigg(
          {{\partial^2 a_y}\over{\partial z\partial x}}
             -{{\partial^2 a_x}\over{\partial z\partial y}}
        \bigg)\cr
      &=\bigg(
          {{\partial^2}\over{\partial y\partial z}}
             -{{\partial^2}\over{\partial z\partial y}}
        \bigg)a_x
       +\bigg(
          {{\partial^2}\over{\partial z\partial x}}
             -{{\partial^2}\over{\partial x\partial z}}
        \bigg)a_y
       +\bigg(
          {{\partial^2}\over{\partial x\partial y}}
             -{{\partial^2}\over{\partial y\partial x}}
        \bigg)a_z\cr
      &=0.\cr
  }
$$
Med andra ord, n{\"a}r vi utvecklar parenteserna f{\"o}rekommer varje blandad
andraderivata tv{\aa} g{\aa}nger med motsatt tecken, vilket via kancellering
resulterar i att divergensen av en rotation alltid {\"a}r identiskt noll.

Den fysikaliska tolkningen {\"a}r att ett f{\"a}lt som representeras som
rotationen av ett ``tillr{\"a}ckligt glatt'' vektorf{\"a}lt automatiskt {\"a}r
divergensfritt, vilket ocks{\aa} {\"a}r precis vad vi ser f{\"o}r, till
exempel, det magnetiska f{\"a}ltet, d{\"a}r vi ser att Gauss lag,
$$
  \nabla\cdot{\bf B}=0,
$$
ger att vi alltid m{\aa}ste kunna formulera f{\"a}ltet f{\"o}r ${\bf B}$ som en
rotation av en vektorfuktion, den s{\aa} kallade vektorpotentialen ${\bf A}$,
som
$$
  {\bf B}=\nabla\times{\bf A}.
$$
Vektoridentiteten $\nabla\cdot(\nabla\times{\bf a})=0$ s{\"a}ger oss
allts{\aa} att vi alltid representera magnetf{\"a}ltet ${\bf B}$ med en
vektorpotential ${\bf A}$ p{\aa} ett s{\aa}dant s{\"a}tt att magnetf{\"a}ltet
per automatik uppfyller Gauss lag f{\"o}r magnetism.
Detta {\"a}r inte en slumpm{\"a}ssig algebraisk f{\"o}renkling, utan {\"a}r en
direkt koppling mellan den matematiska strukturen hos rotationen och den
fysikaliska egenskapen att magnetf{\"a}ltet saknar isolerade k{\"a}llor och
s{\"a}nkor.\numberedfootnote{Att magnetf{\"a}ltet saknar isolerade k{\"a}llor
  och s{\"a}nkor {\"a}r vad vi genomg{\aa}ende i kursen kommer att iterera
  som faktumet att {\it magnetiska monopoler inte existerar}.}

Identiteten har dock samma krav p{\aa} regularitet som tidigare, och om
f{\"a}ltet har singulariteter eller inte {\"a}r ``tillr{\"a}ckligt glatt'',
s{\aa} m{\aa}ste man behandla dessa noggrant, exempelvis med distributioner.

\section{{\hl{Tolkningen av $\nabla\times(\nabla\times{\bf a})$}}}
Den tredje identiteten skiljer sig fr{\aa}n de tv{\aa} f{\"o}reg{\aa}ende,
i det att den inte producerar ett resultat som {\"a}r identiskt noll, utan
ist{\"a}llet formulerar en anv{\"a}ndbar relation mellan differentialoperatorer.
D{\aa} rotationsoperatorn $\nabla\times$ appliceras tv{\aa} g{\aa}nger p{\aa}
ett vektorf{\"a}lt ${\bf a}$, s{\aa} erh{\aa}ller vi
identiteten\numberedfootnote{Notera att vi h{\"a}r v{\"a}ljer att explicit
  beh{\aa}lla parentesen i ``$\nabla\times(\nabla\times{\bf a})$'' f{\"o}r
  de tv{\aa} p{\aa} varandra f{\"o}ljande rotationerna. Eftersom rotationen
  {\"a}r en operator som verkar p{\aa} det som st{\aa}r omedelbart till
  h{\"o}ger, s{\aa} utf{\"o}rs operationerna fr{\aa}n h{\"o}ger till
  v{\"a}nster, det vill s{\"a}ga att vi implicit alltid ser dessa
  operatorer som verkandes enligt
  $$
    \nabla\times\underbrace{(\nabla\times{\bf a})}_{\hbox{1:a rot.}}.
  $$
  Detta {\"a}r precis samma konvention som vi anv{\"a}nder f{\"o}r sammansatta
  funktioner och operatorer i matematiken, exempelvis ``$f(g(x))$'' utan att
  vi inf{\"o}ra n{\aa}gon ytterligare notation varje g{\aa}ng en funktion
  verkar p{\aa} resultatet av en annan funktion. I vektoranalys {\"a}r det
  emellertid en vedertagen praxis att vi lika g{\"a}rna gan skriva detta som
  $$
    \nabla\times\underbrace{\nabla\times{\bf a}}_{\hbox{1:a rot.}}.
  $$
  utan n{\aa}gra explicita parenteser.
  Det {\"a}r allts{\aa} inte n{\aa}gon s{\"a}rskild algebraisk egenskap hos
  rotationen som g{\"o}r parenteserna {\"o}verfl{\"o}diga, utan snarare att
  vi har en etablerad operatornotation och konvention att operatorerna verkar
  fr{\aa}n h{\"o}ger till v{\"a}nster.
  Notera att
  $$
    \nabla\times(\nabla\times{\bf a})\ne(\nabla\times\nabla)\times{\bf a},
  $$
  d{\"a}r uttrycket i h{\"o}gerledet saknar vedertagen betydelse. Det {\"a}r
  h{\"a}r viktigt att inse att $\nabla$ {\it inte} {\"a}r en vanlig vektor,
  utan en {\it differentialoperator}.}
$$
  \nabla\times(\nabla\times{\bf a})
    =\nabla(\nabla\cdot{\bf a})-\nabla^2{\bf a}.
$$\sidx{Vektoridentitet}[$\nabla\times(\nabla\times{\bf a})
    =\nabla(\nabla\cdot{\bf a})-\nabla^2{\bf a}$]
H{\"a}r {\"a}r den skal{\"a}ra operatorn
$$
  \nabla^2
    ={{\partial^2}\over{\partial x^2}}
       +{{\partial^2}\over{\partial y^2}}
       +{{\partial^2}\over{\partial z^2}}.
$$
den s{\aa} kallade {\it Laplaceoperatorn}, vilken vi kommer att studera i
detalj i F{\"o}rel{\"a}sningarna~2 och~3, samt {\aa}terkomma till i
F{\"o}rel{\"a}sning~9 n{\"a}r vi kommer in p{\aa} den elektromagnetiska
v{\aa}gekvationen.\numberedfootnote{Ut{\"o}ver detta s{\aa} {\"a}r
  operatorn $\nabla\times\nabla\times$ central i den f{\"o}rsta
  tentamensuppgiften! Denna tentamensuppgift finns tillsammans med
  ett l{\"o}sningsf{\"o}rslag att h{\"a}mta p{\aa} GitHub-repositoriet
  f{\"o}r kursen, p{\aa}
  {\tt https://github.com/hp35/elmagii/tree/main/examprob}}
D{\aa} den skal{\"a}ra Laplaceoperatorn verkar p{\aa} ett vektorf{\"a}lt
till{\"a}mpas den komponentvis
$$
  \nabla^2{\bf a}={\bf e}_x\nabla^2a_x+{\bf e}_y\nabla^2a_y+{\bf e}_z\nabla^2a_z.
$$
Ett s{\"a}tt att geometriskt tolka uttrycket $\nabla\times\nabla\times{\bf a}$,
som till skillnad fr{\aa}n de tv{\aa} tidigare identiteterna allts{\aa} inte
generellt ger resultatet noll, {\"a}r att betrakta det som en {\it operation
som m{\"a}ter hur det lokala rotationsf{\"a}ltet varierar i rummet}.
Det finns dock en djupare geometrisk tolkning {\"a}n att bara s{\"a}ga att vi
tar rotationen tv{\aa} g{\aa}nger.

\subsection{{\hl{Tolkning 1: ``Cirkulationen av en cirkulation''}}}
Geometriskt inneb{\"a}r uttrycket  $\nabla\times\nabla\times{\bf a}$ att vi
unders{\"o}ker hur rotationsf{\"a}ltet $\nabla\times{\bf a}$ cirkulerar kring
infinitesimala slutna kurvor.
Med andra ord,
\quote{%
  \item{1.}{Den f{\"o}rsta rotationen m{\"a}ter cirkulationen hos det
    ursprungliga f{\"a}ltet ${\bf a}$.}
\item{2.}{Den andra rotationen m{\"a}ter i sin tur cirkulationen hos det
    f{\"o}rsta rotationsf{\"a}ltet $\nabla\times{\bf a}$. Denna andra
    rotationen s{\"a}ger inte bara n{\aa}got om hur mycket det f{\"o}rsta
    rotationsf{\"a}ltet f{\"o}r{\"a}ndras, utan {\"a}ven hur dess variation
    {\"a}r orienterad i rummet.}
}
Fr{\aa}gan {\"a}r dock om vi blir s{\aa} v{\"a}rst mycket klokare av vad detta
{\"a}r, d{\aa} ``cirkulationen av en cirkulation'' {\aa}tminstone f{\"o}r
f{\"o}rfattaren och f{\"o}rel{\"a}saren inte {\"a}r n{\aa}got som ger en
direkt bild av vad som f{\"o}rsigg{\aa}r.

\subsection{{\hl{Tolkning 2: Variation i k{\"a}llstyrka samt utj{\"a}mning
  i rummet}}}
Vi kan g{\"o}ra en alternativ tolkning utifr{\aa}n identiteten
$$
  \nabla\times(\nabla\times{\bf a})=\nabla(\nabla\cdot{\bf a})-\nabla^2{\bf a},
$$
vilken kan f{\"o}rst{\aa}s som en uppdelning i tv{\aa} geometriskt olika bidrag,
$$
  \underbrace{
    \nabla(\nabla\cdot{\bf a})
  }_{
    \vbox{
      \hbox{variation i}
      \hbox{k{\"a}llstyrka}
    }
  }
  -
  \underbrace{
    \vphantom{\nabla(\nabla\cdot{\bf a})}
    \nabla^2{\bf a}
  }_{
    \vbox{
      \hbox{utj{\"a}mning}
      \hbox{av f{\"a}ltet}
    }
  }
$$
Det f{\"o}rsta bidraget kommer fr{\aa}n hur f{\"a}ltets divergens varierar i
rummet, n{\"a}rmare best{\"a}mt {\it gradienten av det skal{\"a}rf{\"a}lt som
utg{\"o}rs av k{\"a}llor}. Om vi till exempel betraktar divergensen av det
elektriska f{\"a}ltet, s{\aa} kan vi se denna divergens som ``k{\"a}llan''.

Det andra bidraget beskriver hur varje komponent av f{\"a}ltet skiljer sig
fr{\aa}n sitt lokala medel\-v{\"a}rde,\numberedfootnote{I F{\"o}rel{\"a}sning~2
  och~3 kommer vi successivt att komma in p{\aa} hur Laplace-operatorn har
  en central betydelse i studiet av elektrostatiska problem, i och med att
  denna operator ligger som grund f{\"o}r {\it Laplace differentialekvation
  f{\"o}r den skal{\"a}ra potentialen}. Vi kommer i ramverket f{\"o}r denna
  ekvation in p{\aa} de medelv{\"a}rdesbildande egenskaperna hos $\nabla^2$
  i och med att vi kommer att h{\"a}rleda {\it entydighetsteoremet f{\"o}r
  den elektrostatiska potentialen}.}
genom Laplace-operatorn $\nabla^2$.
I korthet, s{\aa} kan vi tolka Laplace-operatorn som en ``andra-derivata i
rummet'', d{\"a}r v{\"a}rdet p{\aa} operatorn n{\"a}r den verkar p{\aa} en
funktion beskriver hur funtionsv{\"a}rdet i en omedelbar omgivning av en punkt
{\"a}r, enligt
\quote{%
  \item{1.}{$\nabla^2 f({\bf x})>0\quad\rightarrow\quad$ medelv{\"a}rdet i
    den n{\"a}rmaste omgivningen {\"a}r st{\"o}rre {\"a}n funktionsv{\"a}rdet
    i centrum, och punkten befinner sig i n{\aa}gon mening i en ``lokal dal''.}
  \item{2.}{$\nabla^2 f({\bf x})<0\quad\rightarrow\quad$ medelv{\"a}rdet i
    omgivningen {\"a}r mindre {\"a}n funktionsv{\"a}rdet i centrum, och
    punkten befinner sig i n{\aa}gon mening p{\aa} en ``lokal kulle''.}
  \item{3.}{$\nabla^2 f({\bf x})=0\quad\rightarrow\quad$ medelv{\"a}rdet
    i den n{\"a}rmaste omgivningen sammanfaller, i gr{\"a}nsen, med
    funktionsv{\"a}rdet i centrum.}
}
Detta bist{\aa}r oss med observationen att om f{\"a}ltet {\"a}r divergensfritt,
vilket det exempelvis {\"a}r vid frirymdspropagation av elektromagnetiska
v{\aa}gor eller i {\"o}vrigt laddningsfria omr{\aa}den ({\"a}ven i omgivningar
utanf{\"o}r direkta str{\"o}mmar), det vill s{\"a}ga
$$
  \nabla\cdot{\bf a}=0,
$$
s{\aa} erh{\aa}ller vi att
$$
  \nabla\times(\nabla\times{\bf a})=-\nabla^2{\bf a}.
$$
{\it I k{\"a}llfria problem kan vi d{\"a}rf{\"o}r tolka den dubbla rotationen
som den negativa Laplaceoperatorn av vektorf{\"a}ltet}, vilket {\"a}r en
synnerligen anv{\"a}ndbar omtolkning d{\aa} detta  g{\"o}r att vi kan omvandla
differentialekvationer med dubbla rotationer till uttryck som {\"a}r
signifikant enklare att hantera.
Ett tydligt exempel p{\aa} detta utnyttjande av vektoridentiteten f{\"o}r den
dubbla rotationen {\"a}r den elektromagnetiska v{\aa}gekvationen som vi g{\aa}r
igenom i F{\"o}rel{\"a}sning~9.
\vfill\eject

\section{Sammanfattning av F{\"o}rel{\"a}sning~0 -- Grundl{\"a}ggande
  vektoranalys inom f{\"a}ltteori}
\item{$\bullet$}{Ett {\it f{\"a}lt} {\"a}r en storhet som tilldelar ett
  v{\"a}rde till varje punkt i ett omr{\aa}de av rummet, och eventuellt
  ocks{\aa} till varje tidpunkt. {\it Skal{\"a}ra f{\"a}lt} tilldelar
  varje punkt ett enda tal, s{\"a}g $\phi=\phi({\bf x},t)$, medan
  {\it vektorf{\"a}lt} ist{\"a}llet tilldelar varje punkt en vektor,
  s{\"a}g ${\bf A}={\bf A}({\bf x},t)$.}
\item{$\bullet$}{I {\it elektrostatiska} problem uttrycks det elektriska
  f{\"a}ltet i den {\it skal{\"a}ra potentialen} $\phi$ som
  ${\bf E}=-\nabla\phi$, medan statiska magnetf{\"a}lt har motsvarande
  relation till {\it vektorpotentialen} ${\bf A}$ som
  ${\bf B}=\nabla\times{\bf A}$.}
\item{$\bullet$}{Diracs delta-distribution har i en dimension
  {\it urvalsegenskapen}
  $$
    \delta(x)=0,\ x\ne0,\quad
    \int^{\infty}_{-\infty}\delta(x)\,dx=1,\quad\Rightarrow\quad
    \int^{\infty}_{-\infty}f(x)\delta(x-a)\,dx=f(a),
  $$
  och motsvarande i tre dimensioner
  $$
    \delta^{((3))}({\bf x})=0,\ {\bf x}\ne{\bf x},\quad
    \iiint_{{\Bbb R}^3}\delta^{(3)}({\bf x})\,dV=1,\quad\Rightarrow\quad
    \iiint_{{\Bbb R}^3}f({\bf x})\delta^{(3)}({\bf x}-{\bf x}')\,dV
      =f({\bf x}').
  $$}
\item{$\bullet$}{Gradienten $\nabla\phi({\bf x})$ beskriver
  {\it f{\"o}r{\"a}ndring av ett skal{\"a}rt f{\"a}lt} $\phi({\bf x})$,
  och talar om i vilken riktning det skal{\"a}ra f{\"a}ltet {\"o}kar
  snabbast. I elektrostatik ger den elektriska potentialens negativa
  gradient det elektriska f{\"a}ltet ${\bf E}=-\nabla\phi$.}
\item{$\bullet$}{Divergensen $\nabla\cdot{\bf a}({\bf x})$ {\it beskriver
  k{\"a}llor och s{\"a}nkor} f{\"o}r ett vektorf{\"a}lt ${\bf a}(x)$.
  Elektrisk laddning ger divergens hos det elektriska f{\"a}ltet enligt
  Gauss lag $\nabla\cdot{\bf E}=\rho/\varepsilon_0$, medan det magnetiska
  f{\"a}ltet alltid {\"a}r divergensfritt i den klassiska elektromagnetiska
  teorin, med $\nabla\cdot{\bf B}=0$.}
\item{$\bullet$}{Rotationen $\nabla\times{\bf a}({\bf x})$ beskriver lokal
  cirkulation f{\"o}r ett vektorf{\"a}lt ${\bf a}(x)$. Elektro{\it statiska}
  f{\"a}lt {\"a}r rotationsfria, men ett tidsvarierande magnetf{\"a}lt kan
  ge ett elektriskt f{\"a}lt med nollskild rotation. Magnetf{\"a}ltet har en
  nollskild rotation d{\"a}r det finns elektrisk str{\"o}m i magnetostatik.}
\item{$\bullet$}{{\it Konservativa f{\"a}lt} har potentialer, vilket i sin
  tur inneb{\"a}r att linjeintegralen av ett f{\"a}lt mellan tv{\aa} punkter
  {\it endast beror p{\aa} dess {\"a}ndpunkter}. Konservativa f{\"a}lt och
  dess potentialer bist{\aa}r med ett kraftfullt teoretiskt verktyg inom
  elektromagnetism.}
\item{$\bullet$}{Tre mycket anv{\"a}ndbara vektoridentieter med b{\"a}ring
  p{\aa} elektromagnetism {\"a}r
  $$
    \eqalign{
      &\nabla\times(\nabla\phi)=0,\cr
      &\nabla\cdot(\nabla\times{\bf a})={\bf 0},\cr
      &\nabla\times(\nabla\times{\bf a})
         =\nabla(\nabla\cdot{\bf a})-\nabla^2{\bf a}.\cr
    }
  $$
  \sidx{Vektoridentitet}[$\nabla\times(\nabla\phi)={\bf 0}$]
  \sidx{Vektoridentitet}[$\nabla\cdot(\nabla\times{\bf a})=0$]
  \sidx{Vektoridentitet}[$\nabla\times(\nabla\times{\bf a})
      =\nabla(\nabla\cdot{\bf a})-\nabla^2{\bf a}$]
\item{$\bullet$}{Vektoranalysens identiteter har praktisk och fysikalisk
  betydelse, d{\aa} de hj{\"a}lper oss att f{\"o}rst{\aa} hur potentialer
  uppst{\aa}r och fungerar, varf{\"o}r exempelvis magnetf{\"a}ltet {\"a}r
  divergensfritt och hur elektromagnetiska v{\aa}gor uppst{\aa}r.}
\item{$\bullet$}{Vektoranalys {\"a}r inte bara ett s{\"a}tt att r{\"a}kna
  med vektorer, utan {\"a}r v{\aa}rt ``spr{\aa}k'' f{\"o}r att beskriva
  hur f{\"a}lt varierar, samverkar och propagerar.}
\index
\bye
